\documentclass[UTF8,11pt,a4paper]{ctexart}
\usepackage{amsmath,amssymb,geometry,enumitem,hyperref,array,booktabs,longtable}
\geometry{margin=2cm}\setlist{itemsep=2pt,topsep=4pt}
\raggedbottom
\setlength{\emergencystretch}{2em}
\hypersetup{colorlinks=true,linkcolor=blue}
\newcommand{\R}{\mathbb{R}}
\newcommand{\N}{\mathbb{N}}
\newcommand{\Z}{\mathbb{Z}}
\newcommand{\Q}{\mathbb{Q}}
\newcommand{\BOTH}[2]{\par\noindent\textbf{中文。}\ #1\par\noindent\textbf{English.}\ #2\par}
\newcommand{\EX}[2]{\par\noindent\textbf{例子。}\ #1\par\noindent\textbf{Example.}\ #2\par}
\newcommand{\READ}[1]{\par\noindent\textbf{公式英文读法 / How to read it.} #1\par}
\newcommand{\FORMULA}[2]{\par\noindent\textbf{公式名称 / Formula name.} #1 / #2\par}
\newcommand{\ASK}[4]{\begin{samepage}\par\noindent\textbf{可能问题。}\ #1\par\noindent\textbf{Possible question.}\ #2\par\noindent\textbf{回答要点。}\ #3\par\noindent\textbf{Key answer.}\ #4\medskip\end{samepage}}
\newcommand{\REMARK}[2]{\par\noindent\textbf{注记（应用）。}\ #1\par\noindent\textbf{Remark (application).}\ #2\par}
\newcommand{\APS}[2]{\par\noindent\textbf{中文。} #1\par\noindent\textbf{English.} #2\par\smallskip}
\title{数学分析（合并·完整）APS 面试准备\\[4pt]\large 实数完备性、极限与连续、微积分、级数、含参积分与 Fourier 级数}
\author{Mathematical Analysis (Consolidated): Bilingual Oral Review}
\date{}
\begin{document}\maketitle

\section*{课程概述 | Course overview}
\APS{本讲义把三个学期的数学分析合并为一条严谨主线：实数与完备性、极限与连续、一元微分（中值定理、Taylor）、一元积分（Riemann 可积、微积分基本定理、反常积分）、数项级数（判别法）、多元极限与连续（紧致性）、多元微分（全微分、链式法则、隐函数与反函数、极值与 Lagrange 乘子）、多元积分（Fubini、变量替换、两类线面积分、Green--Gauss--Stokes）、含参积分（连续性、可导性与一致收敛）、Fourier 级数。每个定理都给出完整假设与结论，并配证明思路；每章只保留一两个短例或反例，三册重复的数值例题已删除，应用压缩为注记。}{This text consolidates three semesters into one rigorous main line: the reals and completeness, limits and continuity, one-variable differentiation (mean value theorem, Taylor), one-variable integration (Riemann integrability, the fundamental theorem, improper integrals), numerical series (convergence tests), multivariable limits and continuity (compactness), multivariable differentiation (total differential, chain rule, implicit and inverse function theorems, extrema and Lagrange multipliers), multivariable integration (Fubini, change of variables, line and surface integrals, Green--Gauss--Stokes), integrals depending on a parameter (continuity, differentiation and uniform convergence), and Fourier series. Each theorem states its full hypotheses and conclusion with a proof idea; each chapter keeps only one or two short examples or counterexamples, duplicated numerical examples are removed, and applications are compressed into remarks.}
\APS{整门课可概括为两个问题：局部能否线性化（微分），整体能否由边界控制（积分公式）。面试策略：对每个定理能说出“条件—结论—证明思路—一个反例”，对每个公式说出适用条件，术语用中英双语准确表述。}{The whole course answers two questions: when can a map be linearized locally (differentiation), and when is a global quantity controlled by the boundary (the integral theorems). Interview strategy: for each theorem state hypotheses--conclusion--proof idea--a counterexample, for each formula its conditions of validity, and express terminology accurately in both languages.}

\section*{记号 | Notation}
\BOTH{$\N,\Z,\Q,\R$ 分别是正整数、整数、有理数、实数；$C[a,b]$ 是 $[a,b]$ 上连续函数全体，$C^k$ 是直到 $k$ 阶偏导都连续的函数类。$\sup E$、$\inf E$ 是 $E$ 的上、下确界。$\R^n$ 中点记 $x=(x_1,\ldots,x_n)$，范数 $|x|=\sqrt{x_1^2+\cdots+x_n^2}$，开球 $B(a,r)=\{x:|x-a|<r\}$。区域是连通开集；$\nabla f$、$H_f$、$\nabla\cdot F$、$\nabla\times F$ 分别是梯度、Hessian、散度、旋度。}{$\N,\Z,\Q,\R$ denote the positive integers, integers, rationals and reals; $C[a,b]$ is the space of continuous functions on $[a,b]$, and $C^k$ the class with continuous partials up to order $k$. $\sup E$, $\inf E$ are the supremum and infimum of $E$. In $\R^n$ write $x=(x_1,\ldots,x_n)$ with norm $|x|$ and open ball $B(a,r)$. A region is a connected open set; $\nabla f$, $H_f$, $\nabla\cdot F$, $\nabla\times F$ denote gradient, Hessian, divergence and curl.}

\section{实数与完备性 | The reals and completeness}

\subsection{实数公理与确界原理 | The axioms and the supremum principle}
\BOTH{\textbf{定义（上界与上确界）。}设 $A\subset\R$ 非空。$M$ 是 $A$ 的上界，若对每个 $x\in A$ 有 $x\le M$；此时 $A$ 上有界。若 $A$ 非空且有上界，则其最小上界存在且唯一，称为上确界 $\sup A$。$L=\sup A$ 当且仅当 (1) $L$ 是 $A$ 的上界；(2) 对每个 $\varepsilon>0$ 存在 $x\in A$ 使 $x>L-\varepsilon$。第二条说 $L$ 不能再减小：任何比 $L$ 小的数都不是上界。下确界 $\inf A$（最大下界）对偶地定义。}{\textbf{Definition (upper bound and supremum).} Let $\emptyset\ne A\subset\R$. A number $M$ is an upper bound of $A$ if $x\le M$ for all $x\in A$; then $A$ is bounded above. If $A$ is nonempty and bounded above, its least upper bound exists and is unique, the supremum $\sup A$. $L=\sup A$ iff (1) $L$ is an upper bound of $A$, and (2) for every $\varepsilon>0$ some $x\in A$ exceeds $L-\varepsilon$. The second condition says $L$ cannot be lowered: no smaller number is an upper bound. The infimum $\inf A$ (greatest lower bound) is defined dually.}
\FORMULA{上确界的 $\varepsilon$-刻画}{the epsilon-characterisation of the supremum}
\[
L=\sup A\quad\Longleftrightarrow\quad
\Bigl(\forall x\in A:\ x\le L\Bigr)\ \wedge\
\Bigl(\forall \varepsilon>0\ \exists x\in A:\ x>L-\varepsilon\Bigr).
\]
\BOTH{实数公理分为三组：域公理（四则）、序公理（全序与序相容）和完备性公理（确界原理）。前两组 $\Q$ 也满足，但 $\Q$ 有“洞”。确界原理是实数完备性的公理形式，它是分析中一切存在性定理（极限、零点、最值、反函数、积分）的最终来源。}{The axioms of the reals split into the field axioms (arithmetic), the order axioms (total order compatible with the field), and the completeness axiom (the supremum principle). The rationals satisfy the first two but have gaps. The supremum principle is the axiomatic form of completeness and the ultimate source of every existence theorem in analysis (limits, zeros, extrema, inverse functions, integrals).}
\BOTH{\textbf{推论（Archimedean 性质）。}$\N$ 无上界；故对每个 $\varepsilon>0$ 存在 $n\in\N$ 使 $1/n<\varepsilon$。证：若 $\N$ 有上界，令 $L=\sup\N$，由确界定义存在 $n\in\N$ 使 $n>L-1$，于是 $n+1>L$，与 $L$ 为上界矛盾。这条性质是 $1/n\to0$ 以及一切“取足够大的 $n$”论证的依据，而它本身来自完备性。}{\textbf{Corollary (Archimedean property).} $\N$ is unbounded above; hence for every $\varepsilon>0$ there is $n\in\N$ with $1/n<\varepsilon$. Proof: if $\N$ were bounded above, let $L=\sup\N$; the definition gives $n\in\N$ with $n>L-1$, so $n+1>L$, contradicting that $L$ is an upper bound. This property underlies the limit $1/n\to0$ and every “take $n$ large” argument, and it follows from completeness.}
\EX{$\sup\{1-1/n:n\ge1\}=1$：$1$ 是上界，且给定 $\varepsilon>0$ 取 $n>1/\varepsilon$ 得 $1-\varepsilon<1-1/n<1$；但 $1$ 不属于集合。$\inf\{1/n:n\ge1\}=0$。$\sup(0,1)=1$ 而 $(0,1)$ 无最大值，说明完备性只保证上确界存在，不保证最大值被取到。}{$\sup\{1-1/n:n\ge1\}=1$: one is an upper bound, and given $\varepsilon>0$ the choice $n>1/\varepsilon$ gives $1-\varepsilon<1-1/n<1$; yet $1$ is not in the set. Also $\inf\{1/n:n\ge1\}=0$. $\sup(0,1)=1$ while $(0,1)$ has no maximum, showing completeness guarantees a supremum but not an attained maximum.}

\subsection{有理数不完备：为什么需要实数 | Why the rationals are incomplete}
\BOTH{集合
\[
A=\{q\in\Q:\ q\ge0,\ q^2<2\}
\]
在 $\Q$ 中非空（$0\in A$）、有上界（$2$），但在 $\Q$ 中没有上确界。证分三情形：若 $s^2<2$，取充分大的正整数 $n$ 使 $(2s+1)/n<2-s^2$，则 $(s+1/n)^2=s^2+2s/n+1/n^2<s^2+(2s+1)/n<2$，故 $s+1/n\in A$ 且大于 $s$，与 $s$ 是上界矛盾；若 $s^2>2$，取 $n$ 使 $2s/n<s^2-2$，则 $(s-1/n)^2>s^2-2s/n>2$，而每个 $q\in A$ 满足 $q^2<2$，故 $q<s-1/n$，即 $s-1/n$ 是更小的上界，矛盾；$s^2=2$ 在 $\Q$ 无解（既约 $p/q$ 推出 $p,q$ 皆偶，矛盾）。在 $\R$ 中该上确界就是 $\sqrt2$。}{The set displayed above is nonempty in $\Q$ (since $0\in A$) and bounded above by $2$, but has no supremum in $\Q$. Three cases: if $s^2<2$, choose $n$ with $(2s+1)/n<2-s^2$, then $(s+1/n)^2<2$, so $s+1/n\in A$ exceeds $s$, a contradiction; if $s^2>2$, choose $n$ with $2s/n<s^2-2$, then $(s-1/n)^2>2$, and every $q\in A$ has $q^2<2$, so $q<s-1/n$, a smaller upper bound; $s^2=2$ has no rational solution (a fraction in lowest terms forces both parts even). In $\R$ the supremum is $\sqrt2$.}
\BOTH{用有理端点对 $\sqrt2$ 二分：$a_1=1,b_1=2$，每步取中点 $m_n$，按 $m_n^2<2$ 与否缩小区间，得长度 $2^{1-n}\to0$ 的闭区间套，在 $\R$ 中唯一公共点是 $\sqrt2$，在 $\Q$ 中无公共点。故 Cantor 闭区间套定理在 $\Q$ 失效；$f(x)=x^2-2$ 在 $\Q\cap[1,2]$ 连续、$f(1)=-1<0<2=f(2)$，但无有理零点，介值定理也失效。这两个反例说明“洞”会破坏分析的基本工具。}{Bisect around $\sqrt2$ with rational endpoints: $a_1=1$, $b_1=2$; at each step keep the half containing $\sqrt2$ by testing $m_n^2<2$. This gives nested intervals of length $2^{1-n}\to0$ whose unique common point in $\R$ is $\sqrt2$, with no common point in $\Q$; hence the nested interval theorem fails over $\Q$. The function $f(x)=x^2-2$ is continuous on $\Q\cap[1,2]$ with $f(1)=-1<0<2=f(2)$, yet has no rational zero, so the intermediate value theorem also fails.}

\subsection{完备性的六个等价命题 | Six equivalent formulations of completeness}
\BOTH{在实数系中，下面六个命题等价；它们从集合、数列、区间和覆盖四个角度表达同一条“无缺口”性质。}{For the real numbers, the following six statements are equivalent; they express the same absence of gaps using sets, sequences, intervals and covers.}
\begin{enumerate}[label=(\arabic*)]
\item \textbf{确界原理 / least-upper-bound property.} 每个非空且有上界的 $A\subset\R$ 都有上确界。 Every nonempty set bounded above has a supremum.
\item \textbf{单调有界数列收敛定理 / monotone convergence.} 单调有界实数列收敛；递增有上界时收敛到其上确界。 A bounded monotone sequence converges; an increasing one bounded above converges to its supremum.
\item \textbf{Cantor 闭区间套定理 / nested interval theorem.} 非空闭区间 $I_1\supseteq I_2\supseteq\cdots$ 且长度趋于零，则交恰为一点。 Nested nonempty closed intervals with lengths tending to zero have exactly one common point.
\item \textbf{Bolzano--Weierstrass 定理.} 有界实数列必有收敛子列。 Every bounded real sequence has a convergent subsequence.
\item \textbf{Cauchy 收敛准则 / Cauchy criterion.} 实数列收敛当且仅当它是 Cauchy 列。 A real sequence converges iff it is Cauchy.
\item \textbf{Heine--Borel 定理.} 闭区间 $[a,b]$ 的每个开覆盖有有限子覆盖；等价地，$\R^n$ 中集合紧致当且仅当它闭且有界。 Every open cover of $[a,b]$ has a finite subcover; equivalently a set in $\R^n$ is compact iff it is closed and bounded.
\end{enumerate}
\BOTH{证明等价只需给出循环 $1\Rightarrow2\Rightarrow3\Rightarrow4\Rightarrow5\Rightarrow6\Rightarrow1$，每一步只用到前一步。}{(1)$\Rightarrow$(2) 令 $L=\sup\{a_n\}$，任给 $\varepsilon>0$ 有 $a_N>L-\varepsilon$，递增性使 $n\ge N$ 时 $L-\varepsilon<a_N\le a_n\le L$，故 $a_n\to L$；递减情形对 $-a_n$ 用同一论证。(2)$\Rightarrow$(3) 设 $I_n=[a_n,b_n]$，$a_n\uparrow a$、$b_n\downarrow b$，由 $a_n\le b_n$ 取极限得 $a\le b$，又 $b-a=\lim(b_n-a_n)=0$，公共点即 $a=b$，唯一性由 $|y-x|\le b_n-a_n\to0$。(3)$\Rightarrow$(4) 二分含无穷多项的一半得闭区间套 $I_k$，公共点 $x$ 即子列极限：依次从 $I_k$ 选 $x_{n_k}$ 使 $|x_{n_k}-x|\le|I_k|\to0$。(4)$\Rightarrow$(5) Cauchy 列有界，取收敛子列 $x_{n_k}\to x$，再用 Cauchy 条件把整列拉回：$|x_n-x|\le|x_n-x_{n_k}|+|x_{n_k}-x|<\varepsilon$。(5)$\Rightarrow$(6) 反设无有限子覆盖，二分保持“无有限子覆盖”得闭区间套，取 $x_k\in I_k$ 为 Cauchy 列，其极限落入某覆盖开集，矛盾。(6)$\Rightarrow$(1) 上界集 $\Sigma$ 无最小元时，每点有邻域整体在 $\Sigma$ 内或外，覆盖 $[x_0,M]$，有限子覆盖沿链推回得 $x_0\in\Sigma$，矛盾。}{The equivalence is proved by the cycle $1\Rightarrow2\Rightarrow\cdots\Rightarrow6\Rightarrow1$, each step using only the previous one: (1)$\Rightarrow$(2) put $L=\sup\{a_n\}$; (2)$\Rightarrow$(3) $a_n\uparrow a$, $b_n\downarrow b$ with $b-a=\lim(b_n-a_n)=0$; (3)$\Rightarrow$(4) bisect keeping a half with infinitely many terms; (4)$\Rightarrow$(5) a Cauchy sequence is bounded, extract a convergent subsequence and pull back; (5)$\Rightarrow$(6) bisect preserving “no finite subcover”, the Cauchy points converge into one cover set; (6)$\Rightarrow$(1) each point has a neighbourhood inside or outside the upper-bound set, and a finite subcover chained back gives a contradiction.}
\BOTH{条件不可削弱：闭区间套的“闭”不能去（$(0,1/n)$ 交为空）；BW 的“有界”不能去（$a_n=n$）；Heine--Borel 在无限维空间失效（平方可和空间的标准基两两距离 $\sqrt2$，无收敛子列）。}{The hypotheses cannot be dropped: “closed” in the nested interval theorem ($(0,1/n)$ has empty intersection), “bounded” in Bolzano--Weierstrass ($a_n=n$), and Heine--Borel fails in infinite dimensions (the standard basis vectors of the square-summable space are pairwise at distance $\sqrt2$).}
\REMARK{Dedekind 分割与 Cauchy 完备化都从 $\Q$ 补齐“洞”得到 $\R$，二者给出同构的有序域，说明“无洞”“每个 Cauchy 列收敛”“每个非空有上界集合有上确界”是同一件事。}{Dedekind cuts and Cauchy completion both fill the gaps of $\Q$ to obtain $\R$, giving isomorphic ordered fields; “no gaps”, “every Cauchy sequence converges” and “every nonempty bounded-above set has a supremum” say the same thing.}
\ASK{上确界和最大值有什么区别？}{What is the difference between a supremum and a maximum?}{最大值若存在必属于集合且等于上确界；上确界可能不被取到。$(0,1)$ 有上确界 $1$ 无最大值，$\{1-1/n\}$ 的上确界 $1$ 也不被取到。完备性只保证上确界存在；要取到最大值还需紧性加连续性。}{A maximum, if it exists, belongs to the set and equals the supremum; the supremum may lie outside. $(0,1)$ has supremum $1$ but no maximum, and the supremum $1$ of $\{1-1/n\}$ is not attained. Completeness guarantees a supremum; compactness plus continuity upgrades it to a maximum.}

\section{极限与连续（一元）| Limits and continuity (one variable)}

\subsection{数列极限：定义、性质与单调有界 | Sequences: definition, properties, monotone convergence}
\FORMULA{数列收敛的 $\varepsilon$--$N$ 定义}{the epsilon-N definition of sequence convergence}
\BOTH{$a_n\to L$ 指对每个 $\varepsilon>0$ 存在正整数 $N$，使 $n\ge N$ 时 $|a_n-L|<\varepsilon$。极限唯一；收敛数列有界（取 $\varepsilon=1$，$n\ge N$ 时 $|a_n|\le|L|+1$，前有限项由最大值控制）。四则运算保持极限（商要求分母极限非零）；夹逼定理：若最终 $a_n\le b_n\le c_n$ 且两侧都趋于 $L$，则 $b_n\to L$。数列不收敛到 $L$ 的否定是：存在 $\varepsilon_0>0$，对每个 $N$ 都有某个 $n\ge N$ 使 $|a_n-L|\ge\varepsilon_0$。}{$a_n\to L$ means that for every $\varepsilon>0$ there is $N$ with $|a_n-L|<\varepsilon$ for all $n\ge N$. Limits are unique and convergent sequences are bounded (take $\varepsilon=1$; the finitely many earlier terms are bounded). Arithmetic preserves limits (quotients need a nonzero limiting denominator); squeeze: if $a_n\le b_n\le c_n$ eventually and both ends tend to $L$, then $b_n\to L$. Failure to converge to $L$ means some $\varepsilon_0>0$ has, for every $N$, an index $n\ge N$ with $|a_n-L|\ge\varepsilon_0$.}
\BOTH{单调有界定理：递增有上界的数列收敛到其上确界（递减有下界同理）。Bolzano--Weierstrass：有界数列有收敛子列。Cauchy 准则：实数列收敛当且仅当它是 Cauchy 列。收敛数列的每个子列收敛到同一极限；两个子列极限不同即可否定收敛（如 $(-1)^n$ 的偶奇子列趋于 $1,-1$）。}{Monotone convergence: a nondecreasing sequence bounded above converges to its supremum (dually for nonincreasing). Bolzano--Weierstrass: a bounded sequence has a convergent subsequence. Cauchy criterion: a real sequence converges iff it is Cauchy. Every subsequence of a convergent sequence has the same limit; two subsequences with different limits disprove convergence (e.g. $(-1)^n$).}
\BOTH{阶与无穷小：$a_n=o(b_n)$ 指 $a_n/b_n\to0$，$a_n=O(b_n)$ 指 $|a_n/b_n|$ 最终有界，$a_n\sim b_n$ 指 $a_n/b_n\to1$。等价无穷小可在乘除中替换，加法中主项抵消时不能逐项替换。Stolz 定理（$\infty/\infty$ 型）：$b_n$ 严格递增趋于 $\infty$ 且 $(a_{n+1}-a_n)/(b_{n+1}-b_n)\to L$，则 $a_n/b_n\to L$；$0/0$ 型改用 $a_n,b_n\to0$、$b_n$ 严格递减。}{Orders: $a_n=o(b_n)$ means $a_n/b_n\to0$, $a_n=O(b_n)$ eventual boundedness, $a_n\sim b_n$ means $a_n/b_n\to1$. Equivalent infinitesimals may replace factors in products and quotients, but not terms where leading parts cancel. Stolz ($\infty/\infty$): if $b_n$ strictly increases to $\infty$ and the difference quotient tends to $L$, then $a_n/b_n\to L$; the $0/0$ form uses $a_n,b_n\to0$ with $b_n$ strictly decreasing.}
\EX{证 $1/n\to0$：取 $N>1/\varepsilon$。递推 $a_1=1$、$a_{n+1}=\sqrt{2+a_n}$：归纳得 $0<a_n<2$ 且 $a_{n+1}^2-a_n^2=(2-a_n)(1+a_n)>0$，故递增有界，收敛到 $L$，由 $\sqrt{\ }$ 连续得 $L=\sqrt{2+L}$，故 $L=2$。}{To prove $1/n\to0$ take $N>1/\varepsilon$. For $a_1=1$, $a_{n+1}=\sqrt{2+a_n}$, induction gives $0<a_n<2$ and $a_{n+1}^2-a_n^2=(2-a_n)(1+a_n)>0$, so the sequence is increasing and bounded, hence converges to $L$ with $L=\sqrt{2+L}$, so $L=2$.}
\ASK{Cauchy 准则和单调有界定理分别适合什么情形？}{When are the Cauchy criterion and monotone convergence useful?}{若能直接比较相邻尾项或尾和（级数、递推估计），Cauchy 准则合适，不需要预先知道候选极限；若能证单调并找到统一上下界，单调有界定理更短，极限往往可由上确界或递推关系得到。两者都依赖实数完备性。}{Use the Cauchy criterion when tail differences or sums can be bounded directly (series, recursion estimates); it needs no candidate limit in advance. Use monotone convergence when monotonicity and a uniform bound are easy; it is shorter and the limit is often read from the supremum or recursion. Both rely on completeness.}

\subsection{函数极限、连续与间断 | Function limits, continuity and discontinuities}
\FORMULA{函数极限的 $\varepsilon$--$\delta$ 定义}{the epsilon-delta definition of a function limit}
\BOTH{设 $a$ 是定义域 $D$ 的聚点。$\lim_{x\to a}f(x)=L$ 指对每个 $\varepsilon>0$ 存在 $\delta>0$，使 $x\in D$、$0<|x-a|<\delta$ 时 $|f(x)-L|<\varepsilon$。$a$ 可以不属于定义域。左右极限分别限制 $x>a$、$x<a$；双侧极限存在当且仅当两者存在且相等。无穷远处把条件换成 $|x|>R$。}{Let $a$ be an accumulation point of the domain $D$. $\lim_{x\to a}f(x)=L$ means for every $\varepsilon>0$ there is $\delta>0$ with $|f(x)-L|<\varepsilon$ whenever $x\in D$ and $0<|x-a|<\delta$; $a$ need not lie in $D$. One-sided limits restrict to $x>a$ or $x<a$; the two-sided limit exists iff both exist and agree. At infinity the condition becomes $|x|>R$.}
\BOTH{\textbf{Heine 定理。}$\lim_{x\to a}f(x)=L$ 当且仅当对每个 $x_n\to a$（$x_n\ne a$）都有 $f(x_n)\to L$。它既能证明极限，也能用两条不同趋近数列否定极限。基本极限：$\sin x/x\to1$，$(1+x)^{1/x}\to e$（$x\to0$）。函数极限也有 Cauchy 准则。}{Heine's theorem: $\lim_{x\to a}f(x)=L$ iff every sequence $x_n\to a$ with $x_n\ne a$ satisfies $f(x_n)\to L$. It proves limits and disproves them via two approaching sequences. Basic limits: $\sin x/x\to1$, $(1+x)^{1/x}\to e$. Function limits also satisfy a Cauchy criterion.}
\BOTH{$f$ 在 $a$ 连续指 $\lim_{x\to a}f(x)=f(a)$；三种等价说法：$\varepsilon$--$\delta$ 判据、序列判据（$x_n\to a\Rightarrow f(x_n)\to f(a)$）、开集原像相对开。间断点分类：第一类（左右极限有限，相等为可去、不等为跳跃）；第二类（至少一侧极限不存在，无穷或振荡）。单调函数的间断都是跳跃且至多可数。}{$f$ is continuous at $a$ if $\lim_{x\to a}f(x)=f(a)$; three equivalent forms: the $\varepsilon$--$\delta$ criterion, the sequential criterion, and open preimages being relatively open. Discontinuities: first kind (finite one-sided limits; equal gives removable, unequal a jump); second kind (one side fails, infinite or oscillatory). A monotone function has only jumps, at most countably many.}
\EX{证 $\lim_{x\to2}x^2=4$：$|x^2-4|=|x-2||x+2|$，当 $|x-2|<1$ 时 $|x+2|<5$，取 $\delta=\min(1,\varepsilon/5)$。$f(x)=\sin(1/x)$ 在 $0$ 无极限：$x_n=1/(n\pi)$、$y_n=1/(2n\pi+\pi/2)$ 都趋于 $0$ 而 $f(x_n)=0$、$f(y_n)=1$。}{To prove $\lim_{x\to2}x^2=4$: $|x^2-4|=|x-2||x+2|$ and $|x-2|<1$ gives $|x+2|<5$, so $\delta=\min(1,\varepsilon/5)$. The limit of $\sin(1/x)$ at $0$ fails: $x_n=1/(n\pi)$ and $y_n=1/(2n\pi+\pi/2)$ both tend to $0$ while $f(x_n)=0$ and $f(y_n)=1$.}

\subsection{闭区间上连续函数的性质与一致连续 | Properties on closed intervals and uniform continuity}
\BOTH{\textbf{最值定理 / extreme value theorem.} $f\in C[a,b]$ 则 $f$ 有界且取到最大值与最小值。\textbf{介值定理 / intermediate value theorem.} $f\in C[a,b]$ 取遍 $f(a)$ 与 $f(b)$ 之间的一切值，故把区间映成区间。两者都依赖闭区间的紧性：开区间上可能失效，如 $1/x$ 在 $(0,1)$ 无界、$x^2$ 在 $(0,1)$ 取不到上确界。}{\textbf{Extreme value theorem.} A continuous function on $[a,b]$ is bounded and attains its maximum and minimum. \textbf{Intermediate value theorem.} A continuous function on $[a,b]$ takes every value between $f(a)$ and $f(b)$, so intervals map to intervals. Both rely on compactness: they fail on open intervals, e.g. $1/x$ is unbounded on $(0,1)$ and $x^2$ does not attain its supremum there.}
\BOTH{\textbf{一致连续 / uniform continuity.} 对每个 $\varepsilon>0$ 存在 $\delta>0$，使区间内任意 $|x-y|<\delta$ 都有 $|f(x)-f(y)|<\varepsilon$；$\delta$ 只依赖 $\varepsilon$，全区间统一。\textbf{Cantor 定理.} $f\in C[a,b]$ 则一致连续。证：反设存在 $\varepsilon_0>0$ 及 $x_n,y_n$ 使 $|x_n-y_n|<1/n$ 但 $|f(x_n)-f(y_n)|\ge\varepsilon_0$；由 BW 取 $x_n$ 的收敛子列，$y_n$ 同趋于该点，连续性矛盾。连续不推一致连续：$x^2$ 在 $\R$、$1/x$ 在 $(0,1)$、$\sin(1/x)$ 在 $(0,1)$。一致连续保证连续函数 Riemann 可积。}{\textbf{Uniform continuity.} For every $\varepsilon>0$ there is $\delta>0$ such that $|x-y|<\delta$ implies $|f(x)-f(y)|<\varepsilon$, with $\delta$ depending only on $\varepsilon$. \textbf{Cantor's theorem.} A continuous function on $[a,b]$ is uniformly continuous. Proof: otherwise there are $\varepsilon_0>0$ and $x_n,y_n$ with $|x_n-y_n|<1/n$ but $|f(x_n)-f(y_n)|\ge\varepsilon_0$; a convergent subsequence of $x_n$ (Bolzano--Weierstrass) has $y_n$ tending to the same point, contradicting continuity. Continuity need not be uniform: $x^2$ on $\R$, $1/x$ on $(0,1)$, $\sin(1/x)$ on $(0,1)$. Uniform continuity ensures Riemann integrability of continuous functions.}
\EX{若 $f$ 连续、$f(1)=-2$、$f(3)=5$，则存在 $c\in(1,3)$ 使 $f(c)=0$（介值定理只保证存在，不给个数与位置）。$x^2$ 不一致连续：$x_n=\sqrt{n+1}$、$y_n=\sqrt n$ 有 $x_n-y_n\to0$ 但 $|x_n^2-y_n^2|=1$；$1/x$ 不一致连续：$x_n=1/n$、$y_n=1/(n+1)$ 距离趋于零而函数值之差为 $1$。$\sin x$ 一致连续（$|\sin x-\sin y|\le|x-y|$）。}{If $f$ is continuous with $f(1)=-2$, $f(3)=5$, some $c\in(1,3)$ has $f(c)=0$ (the theorem guarantees existence, not number or location). $x^2$ is not uniformly continuous: $x_n=\sqrt{n+1}$, $y_n=\sqrt n$ have $x_n-y_n\to0$ but $|x_n^2-y_n^2|=1$; $1/x$ is not uniformly continuous on $(0,1)$: with $x_n=1/n$, $y_n=1/(n+1)$ the distance tends to zero while the difference is $1$. By contrast $\sin x$ is uniformly continuous.}
\ASK{数列极限和函数极限的区别是什么？}{What distinguishes a sequence limit from a function limit?}{数列沿离散指标，条件较弱；函数极限要求任意方式趋近，更强。Heine 定理连接两者：函数极限存在当且仅当所有趋近数列的像收敛到同一值，故一条数列即可否定函数极限。例：$\sin(\pi x)$ 在整数上恒为零故 $f(n)\to0$，但 $x\to+\infty$ 无极限。}{A sequence limit follows a discrete index and is weaker; a function limit requires approach in every way and is stronger. Heine's theorem links them: the function limit exists iff every approaching sequence has images converging to the same value, so one sequence can disprove it. Example: $\sin(\pi x)$ vanishes on the integers, so $f(n)\to0$, yet it has no limit as $x\to+\infty$.}

\section{一元微分 | Differentiation of one variable}

\subsection{导数、微分与求导规则 | Derivatives, differentials and rules}
\FORMULA{导数的差商定义与一阶线性近似}{the difference-quotient definition and linearization}
\BOTH{$f'(a)=\lim_{h\to0}[f(a+h)-f(a)]/h$ 是切线斜率，也是局部线性近似系数：$f(a+h)=f(a)+f'(a)h+o(h)$；微分 $df=f'(a)h$ 是主线性部分，一元情形可导与可微等价，可导必连续。求导规则：$(fg)'=f'g+fg'$、$(f/g)'=(f'g-fg')/g^2$、链式法则 $(f\circ g)'=f'(g)g'$、反函数 $(f^{-1})'(y_0)=1/f'(x_0)$（$f'(x_0)\ne0$）。}{$f'(a)=\lim_{h\to0}[f(a+h)-f(a)]/h$ is the tangent slope and the coefficient of the linearization $f(a+h)=f(a)+f'(a)h+o(h)$; the differential $df=f'(a)h$ is the principal linear part. In one variable differentiability equals existence of the differential and implies continuity. Rules: product, quotient, chain rule, and the inverse rule ($(f^{-1})'(y_0)=1/f'(x_0)$ when $f'(x_0)\ne0$).}
\BOTH{高阶导数是反复求导；Leibniz 公式 $(fg)^{(n)}=\sum_{k=0}^n\binom nk f^{(k)}g^{(n-k)}$ 是乘积求导的二项式版本。Darboux 定理：导数具有介值性——若 $f$ 区间内可导且 $f'(a)<\lambda<f'(b)$，则存在 $c$ 使 $f'(c)=\lambda$；故导数不能有跳跃间断（可有振荡间断，如 $x^2\sin(1/x)$ 的导数）。可导必连续，连续未必可导：$|x|$ 在 $0$ 左右导数不等；$x\sin(1/x)$（$f(0)=0$）连续但差商 $\sin(1/h)$ 无极限。}{Higher derivatives are iterated; Leibniz's formula is the binomial theorem for products. Darboux's theorem: derivatives have the intermediate value property — if $f$ is differentiable and $f'(a)<\lambda<f'(b)$, some $c$ has $f'(c)=\lambda$; hence a derivative cannot have a jump (though it may have an oscillatory discontinuity). Differentiability implies continuity, not conversely: $|x|$ has unequal one-sided derivatives at $0$; $x\sin(1/x)$ with $f(0)=0$ is continuous but its difference quotient $\sin(1/h)$ has no limit.}
\EX{$f(x)=x^2$：$f(a+h)-f(a)=2ah+h^2$，微分近似 $2ah$、误差 $h^2$。$(xe^x)^{(n)}=(x+n)e^x$（Leibniz 公式只有 $k=0,1$ 两项）。}{For $x^2$, $f(a+h)-f(a)=2ah+h^2$; the differential keeps $2ah$ with error $h^2$. $(xe^x)^{(n)}=(x+n)e^x$ (only $k=0,1$ survive in Leibniz's formula).}

\subsection{中值定理：Fermat、Rolle、Lagrange、Cauchy | The mean value theorems}
\BOTH{\textbf{Fermat 引理.} 内点局部极值且可导则 $f'=0$。证：极小值处左差商 $\le0$、右差商 $\ge0$，令 $h\to0$。\textbf{Rolle 定理.} $f\in C[a,b]$ 在 $(a,b)$ 可导且 $f(a)=f(b)$，则存在 $c\in(a,b)$ 使 $f'(c)=0$。\textbf{Lagrange 中值定理.} $f\in C[a,b]$ 在 $(a,b)$ 可导，则存在 $c$ 使 $f(b)-f(a)=f'(c)(b-a)$；证明对辅助函数 $F(x)=f(x)-f(a)-\frac{f(b)-f(a)}{b-a}(x-a)$ 用 Rolle。\textbf{Cauchy 中值定理.} $f,g\in C[a,b]$ 在 $(a,b)$ 可导且 $g'\ne0$，则存在 $c$ 使 $\frac{f(b)-f(a)}{g(b)-g(a)}=\frac{f'(c)}{g'(c)}$；取 $g(x)=x$ 即 Lagrange。}{\textbf{Fermat's lemma.} An interior local extremum of a differentiable function has $f'=0$; compare one-sided difference quotients. \textbf{Rolle's theorem.} If $f\in C[a,b]$ is differentiable on $(a,b)$ with $f(a)=f(b)$, then $f'(c)=0$ for some $c\in(a,b)$. \textbf{Lagrange's mean value theorem.} If $f\in C[a,b]$ is differentiable on $(a,b)$, some $c$ has $f(b)-f(a)=f'(c)(b-a)$; apply Rolle to the displayed auxiliary function. \textbf{Cauchy's mean value theorem.} If $f,g\in C[a,b]$ are differentiable on $(a,b)$ with $g'\ne0$, some $c$ satisfies the displayed ratio identity; $g(x)=x$ recovers Lagrange.}
\BOTH{Rolle 三条件缺一不可：$f(x)=x$ 于 $[0,1]$（端点不等）；$|x|$ 于 $[-1,1]$（$0$ 不可导）；$f(x)=x$ 于 $(0,1)$ 且 $f(0)=f(1)=0$（端点不连续）。与单调/极值/凸性的联系：$f'\ge0$ 则单调不减（对 $x<y$ 用 Lagrange）、$f'\equiv0$ 为常数、$f'>0$ 严格递增；$f''\ge0$ 为凸（图像在切线上方、弦下方）。}{Each Rolle hypothesis is essential: $f(x)=x$ on $[0,1]$ (endpoints differ); $|x|$ on $[-1,1]$ (not differentiable at $0$); $f(x)=x$ on $(0,1)$ with $f(0)=f(1)=0$ (discontinuous at an endpoint). Connections: $f'\ge0$ gives monotonicity (Lagrange on $x<y$), $f'\equiv0$ constancy, $f'>0$ strict increase; $f''\ge0$ convexity (graph above tangents, below chords).}
\EX{$f(x)=x^2$ 于 $[1,3]$：$8=f'(c)\cdot2=2c\cdot2$ 得 $c=2$。应用：$|\sin x-\sin y|\le|x-y|$；$x>0$ 时 $e^x-1=e^{\xi}x\ge x$。}{For $x^2$ on $[1,3]$, $8=f'(c)\cdot2=2c\cdot2$ gives $c=2$. Applications: $|\sin x-\sin y|\le|x-y|$, and $e^x-1=e^{\xi}x\ge x$ for $x>0$.}
\ASK{Rolle、Lagrange、Cauchy 中值定理之间是什么关系？}{How are Rolle's, Lagrange's and Cauchy's mean value theorems related?}{Rolle 是端点等高的基础情形；Lagrange 用割线构造辅助函数化归 Rolle；Cauchy 用两个函数的辅助函数化归 Rolle，取 $g(x)=x$ 即 Lagrange。三者都要求闭区间连续、开区间可导，Cauchy 还需 $g'\ne0$。}{Rolle is the equal-endpoint base case; Lagrange reduces to Rolle via the secant auxiliary function; Cauchy reduces to Rolle with an auxiliary function of two functions, and $g(x)=x$ recovers Lagrange. All require continuity on the closed interval and differentiability inside; Cauchy also needs $g'\ne0$.}

\subsection{L'H\^opital 法则 | L'H\^opital's rule}
\BOTH{用于 $0/0$ 或 $\infty/\infty$ 型。$0/0$ 型：$f,g$ 在 $(a,b)$ 可导、$g'\ne0$、$\lim_{x\to a+}f=\lim g=0$，若 $\lim f'/g'=A$（有限或无穷），则 $\lim f/g=A$。$\infty/\infty$ 型只需分母趋于 $\infty$。证明用 Cauchy 中值定理。每次重复使用前都要重新检查不定式型。逆命题不成立：$x^2\sin(1/x)/x$ 的原比趋于 $0$，但导数比含 $\cos(1/x)$ 无极限，法则失效时应改夹逼或等价无穷小。}{For $0/0$: if $f,g$ are differentiable on $(a,b)$, $g'\ne0$, $\lim_{x\to a+}f=\lim g=0$ and $\lim f'/g'=A$ (finite or infinite), then $\lim f/g=A$. For $\infty/\infty$ only the denominator need tend to infinity. The proof uses Cauchy's theorem, and the indeterminate form must be rechecked before each repeated use. The converse fails: $x^2\sin(1/x)/x$ has ratio tending to $0$, but the derivative ratio contains $\cos(1/x)$ and has no limit; then use squeezing or equivalent infinitesimals.}
\EX{$(e^x-1)/x$ 为 $0/0$ 型，求导得 $e^x\to1$；$(\sin x-x)/x^3$ 三次使用法则得 $-1/6$；$(\log x)/x$ 为 $\infty/\infty$ 型，得 $1/x\to0$。}{$(e^x-1)/x$ is $0/0$; differentiating gives $e^x\to1$. Three applications to $(\sin x-x)/x^3$ give $-1/6$. $(\log x)/x$ is $\infty/\infty$; differentiating gives $1/x\to0$.}

\subsection{Taylor 公式与余项 | Taylor's formula and remainders}
\FORMULA{Taylor 公式及 Lagrange 余项}{Taylor's theorem with the Lagrange remainder}
\BOTH{若 $f$ 足够光滑，$f(a+h)=\sum_{k=0}^{n}\frac{f^{(k)}(a)}{k!}h^k+R_n$。Lagrange 余项 $R_n=\frac{f^{(n+1)}(\xi)}{(n+1)!}h^{n+1}$（$\xi$ 在 $a$ 与 $a+h$ 间）用于误差估计；Peano 余项 $R_n=o(h^n)$（只需 $n$ 阶可导）用于求极限；Cauchy 余项与积分余项是另两种形式。常用 Maclaurin 展开：$e^x=1+x+x^2/2!+\cdots$，$\sin x=x-x^3/3!+\cdots$，$\cos x=1-x^2/2!+\cdots$，$\log(1+x)=x-x^2/2+x^3/3-\cdots$（$-1<x\le1$），$(1+x)^\alpha=1+\alpha x+\frac{\alpha(\alpha-1)}{2!}x^2+\cdots$。}{If $f$ is smooth enough, $f(a+h)=\sum_{k=0}^{n}\frac{f^{(k)}(a)}{k!}h^k+R_n$. The Lagrange remainder $R_n=\frac{f^{(n+1)}(\xi)}{(n+1)!}h^{n+1}$ (with $\xi$ between $a$ and $a+h$) gives error bounds; the Peano remainder $R_n=o(h^n)$ (needing only $n$-fold differentiability) serves limits; the Cauchy and integral remainders are two further forms. Common Maclaurin series: $e^x$, $\sin x$, $\cos x$, $\log(1+x)$ for $-1<x\le1$, $(1+x)^\alpha$.}
\EX{$e^{0.1}$：二次 Taylor 多项式得 $1+0.1+0.1^2/2=1.105$，Lagrange 余项 $|R_2|\le e^{0.1}0.1^3/6<0.000185$，真值 $1.105170918\ldots$。}{For $e^{0.1}$ the quadratic Taylor polynomial gives $1.105$; the Lagrange remainder satisfies $|R_2|\le e^{0.1}0.1^3/6<0.000185$, against the exact $1.105170918\ldots$.}
\BOTH{Taylor 级数只在余项趋于零时等于函数。反例 $f(x)=e^{-1/x^2}$（$f(0)=0$）所有阶导数在 $0$ 为零，Maclaurin 级数是零级数，但 $x\ne0$ 时 $f>0$，故 $C^\infty$ 不必解析。用 Taylor 判断极值：若 $f'(c)=\cdots=f^{(m-1)}(c)=0$、$f^{(m)}(c)\ne0$，则 $m$ 偶且 $f^{(m)}(c)>0$ 为极小、$<0$ 为极大，$m$ 奇不是极值。}{A Taylor series equals the function only when its remainder tends to zero. For $f(x)=e^{-1/x^2}$ with $f(0)=0$ all derivatives vanish at $0$, so the Maclaurin series is the zero series while $f>0$ off zero; $C^\infty$ need not be analytic. For extrema, if the first $m-1$ derivatives vanish and $f^{(m)}(c)\ne0$, then even $m$ with $f^{(m)}(c)>0$ ($<0$) gives a minimum (maximum), odd $m$ no extremum.}
\REMARK{Newton 迭代 $x_{n+1}=x_n-f(x_n)/f'(x_n)$ 在单根附近二次收敛（误差递推 $e_{n+1}\approx\frac{f''}{2f'}e_n^2$），但初值不合适会发散（$x^{1/3}=0$ 给出 $x_{n+1}=-2x_n$）。近似值必须配误差界，这是 Taylor 余项的核心用途。}{Newton's iteration $x_{n+1}=x_n-f(x_n)/f'(x_n)$ converges quadratically near a simple root (error recurrence $e_{n+1}\approx\frac{f''}{2f'}e_n^2$) but can diverge otherwise ($x^{1/3}=0$ gives $x_{n+1}=-2x_n$). An approximation must come with an error bound, the core use of the Taylor remainder.}
\ASK{Taylor 公式主要用来做什么？}{What is Taylor's formula used for?}{近似计算并给误差界（如 $e^{0.1}$）；用 Peano 余项求极限；判断极值、凸性与拐点。关键是说明展开点、阶数和余项形式，只有余项趋于零才是函数值。}{Approximation with error bounds (e.g. $e^{0.1}$); limits via the Peano remainder; extrema, convexity and inflection points. State the expansion point, order and remainder form; the series equals the function only when the remainder tends to zero.}

\section{一元积分 | Integration of one variable}

\subsection{原函数与积分技巧 | Antiderivatives and techniques}
\BOTH{若 $F'=f$，$F$ 是原函数；不定积分表示全部原函数，常数 $C$ 不能省。基本表：$\int x^\alpha dx=x^{\alpha+1}/(\alpha+1)+C$（$\alpha\ne-1$）、$\int dx/x=\log|x|+C$、$\int e^xdx=e^x+C$、$\int\sin x\,dx=-\cos x+C$、$\int dx/(1+x^2)=\arctan x+C$。换元法逆用链式法则，分部积分 $\int u\,dv=uv-\int v\,du$ 来自乘积求导。}{If $F'=f$, then $F$ is an antiderivative; the indefinite integral denotes all antiderivatives and includes $C$. Basic table: $\int x^\alpha dx$, $\int dx/x$, $\int e^xdx$, $\int\sin x\,dx$, $\int dx/(1+x^2)$. Substitution reverses the chain rule, integration by parts follows from the product rule.}
\EX{$\int2x\cos(x^2)dx=\sin(x^2)+C$（$u=x^2$）；$\int xe^xdx=xe^x-e^x+C$；$\int\log x\,dx=x\log x-x+C$。}{$\int2x\cos(x^2)dx=\sin(x^2)+C$ with $u=x^2$; $\int xe^xdx=xe^x-e^x+C$; $\int\log x\,dx=x\log x-x+C$.}

\subsection{Riemann 积分与可积性 | The Riemann integral and integrability}
\BOTH{设 $f$ 在 $[a,b]$ 有界，分割 $P:a=x_0<\cdots<x_n=b$，$\Delta x_i=x_i-x_{i-1}$，网格 $\|P\|=\max\Delta x_i$，带点 Riemann 和 $\sigma=\sum f(\xi_i)\Delta x_i$。若存在 $I$，使对每个 $\varepsilon>0$ 存在 $\delta>0$，只要 $\|P\|<\delta$（对任意选点）就有 $|\sigma-I|<\varepsilon$，则 $f$ Riemann 可积，$I=\int_a^bf$。}{For bounded $f$ on $[a,b]$, with mesh $\|P\|=\max\Delta x_i$ and tagged sums $\sigma=\sum f(\xi_i)\Delta x_i$, if some $I$ satisfies that every tagged partition with $\|P\|<\delta$ has $|\sigma-I|<\varepsilon$, then $f$ is Riemann integrable with integral $I$.}
\BOTH{\textbf{Darboux 判据.} 记 $m_i=\inf_{\Delta_i}f$、$M_i=\sup_{\Delta_i}f$，下和 $L=\sum m_i\Delta x_i$、上和 $U=\sum M_i\Delta x_i$，振幅 $\omega_i=M_i-m_i$。有界 $f$ 可积当且仅当对每个 $\varepsilon>0$ 存在分割使 $U-L<\varepsilon$，即 $\sum\omega_i\Delta x_i<\varepsilon$。}{\textbf{Darboux criterion.} With lower and upper sums $L=\sum m_i\Delta x_i$, $U=\sum M_i\Delta x_i$ and oscillation $\omega_i=M_i-m_i$, a bounded $f$ is integrable iff for every $\varepsilon>0$ some partition has $U-L<\varepsilon$, i.e. $\sum\omega_i\Delta x_i<\varepsilon$.}
\BOTH{可积函数类：闭区间连续函数（Cantor 定理给一致连续，从而细分割）、单调函数、有限个跳跃的有界函数。Dirichlet 函数不可积（任何分割 $L=0$、$U=b-a$）。性质：线性、区间可加、单调、$|\int f|\le\int|f|$、$m(b-a)\le\int f\le M(b-a)$、积分中值定理（$f$ 连续时存在 $\xi$ 使 $\int_a^bf=f(\xi)(b-a)$）。}{Integrable classes: continuous functions on a closed interval (Cantor's theorem gives uniform continuity, hence fine partitions), monotone functions, and bounded functions with finitely many jumps. The Dirichlet function is not integrable (every partition has $L=0$, $U=b-a$). Properties: linearity, additivity, monotonicity, $|\int f|\le\int|f|$, the inequality $m(b-a)\le\int f\le M(b-a)$, and the integral mean value theorem.}
\EX{$\int_0^1x^2dx=1/3$：右端 Riemann 和 $\frac1n\sum_{k=1}^n(k/n)^2=\frac{(n+1)(2n+1)}{6n^2}\to1/3$，直接体现定义。}{$\int_0^1x^2dx=1/3$: the right-endpoint Riemann sum $\frac1n\sum(k/n)^2=\frac{(n+1)(2n+1)}{6n^2}\to1/3$, directly illustrating the definition.}
\ASK{Riemann 可积是什么意思？哪些函数可积？}{What does Riemann integrability mean and which functions are integrable?}{所有足够细的带点分割的 Riemann 和趋于同一个数，等价于 Darboux 上下和之差可任意小。连续函数、单调函数、有限个跳跃的有界函数可积，Dirichlet 函数不可积。}{All Riemann sums over sufficiently fine tagged partitions approach the same number, equivalently the Darboux upper and lower sums become arbitrarily close. Continuous functions, monotone functions and bounded functions with finitely many jumps are integrable; the Dirichlet function is not.}

\subsection{微积分基本定理 | The fundamental theorem of calculus}
\FORMULA{微积分基本定理与 Newton--Leibniz 公式}{the fundamental theorem of calculus}
\BOTH{\textbf{第一基本定理.} $f$ 可积时 $F(x)=\int_a^xf(t)dt$ 连续；$f$ 在 $x_0$ 连续时 $F'(x_0)=f(x_0)$。证：差商 $\frac1h\int_{x_0}^{x_0+h}[f(t)-f(x_0)]dt$ 的被积函数在 $x_0$ 附近一致小。\textbf{第二基本定理.} $f$ 可积且 $F'=f$（端点取单侧），则 $\int_a^bf=F(b)-F(a)$；证：对分割在每个子区间用 Lagrange 中值定理，把 $F(b)-F(a)$ 写成 Riemann 和 $\sum f(\xi_i)\Delta x_i$，令网格趋于零。}{\textbf{First theorem.} If $f$ is integrable, $F(x)=\int_a^xf(t)dt$ is continuous; if $f$ is continuous at $x_0$, $F'(x_0)=f(x_0)$. Proof: the difference quotient $\frac1h\int_{x_0}^{x_0+h}[f(t)-f(x_0)]dt$ has uniformly small integrand near $x_0$. \textbf{Second theorem.} If $f$ is integrable with $F'=f$ (one-sided at endpoints), then $\int_a^bf=F(b)-F(a)$; proof: Lagrange on each subinterval expresses $F(b)-F(a)$ as the Riemann sum $\sum f(\xi_i)\Delta x_i$, and the mesh tends to zero.}
\FORMULA{变上限积分求导与换元}{differentiation of a variable-limit integral and substitution}
\[
\frac{d}{dx}\int_{\alpha(x)}^{\beta(x)}f(t)\,dt=f(\beta(x))\beta'(x)-f(\alpha(x))\alpha'(x),\qquad
\int_a^b f(g(x))g'(x)\,dx=\int_{g(a)}^{g(b)}f(u)\,du .
\]
\BOTH{换元须同步变换微分与上下限；定积分分部 $\int_a^bu v'dx=[uv]_a^b-\int_a^bu'v\,dx$，端点项不能漏。}{Substitution transforms both the differential and the endpoints; definite integration by parts is $\int_a^bu v'dx=[uv]_a^b-\int_a^bu'v\,dx$, and the endpoint term must not be dropped.}
\EX{$\int_1^2x^2dx=x^3/3\big|_1^2=7/3$；$\frac{d}{dx}\int_0^{x^2}\sin\sqrt t\,dt=2x\sin x$（$x>0$，变上限加链式）；$\int_0^{\pi}x\sin x\,dx=\pi$。}{$\int_1^2x^2dx=7/3$; $\frac{d}{dx}\int_0^{x^2}\sin\sqrt t\,dt=2x\sin x$ for $x>0$ (variable limit plus chain rule); $\int_0^{\pi}x\sin x\,dx=\pi$.}

\subsection{反常积分与收敛判别 | Improper integrals and convergence tests}
\BOTH{无穷积分 $\int_a^\infty f=\lim_{R\to\infty}\int_a^Rf$；瑕积分在端点 $b$ 处为 $\lim_{c\to b-}\int_a^cf$，内部奇点须拆成两侧分别收敛。反常积分也有 Cauchy 准则。$p$-判别法：$\int_1^\infty x^{-p}dx$ 收敛当且仅当 $p>1$，$\int_0^1x^{-p}dx$ 收敛当且仅当 $p<1$。}{Infinite integrals are truncation limits $\int_a^\infty f=\lim_{R\to\infty}\int_a^Rf$; singular integrals at $b$ are $\lim_{c\to b-}\int_a^cf$, with interior singularities split. Improper integrals satisfy a Cauchy criterion. $p$-test: $\int_1^\infty x^{-p}$ converges iff $p>1$, and $\int_0^1x^{-p}$ iff $p<1$.}
\BOTH{比较判别：$0\le f\le g$ 时 $\int g$ 收敛推 $\int f$ 收敛，$\int f$ 发散推 $\int g$ 发散；极限比较在 $f/g\to c\in(0,\infty)$ 时同敛散。绝对收敛推出收敛。Dirichlet 判别：$F(R)=\int_a^Rf$ 有界、$g$ 单调趋于 $0$，则 $\int_a^\infty fg$ 收敛；Abel 判别：$\int_a^\infty f$ 收敛、$g$ 单调有界，则 $\int_a^\infty fg$ 收敛。}{Comparison: if $0\le f\le g$, convergence of $\int g$ implies that of $\int f$, and divergence of $\int f$ implies that of $\int g$; limit comparison with $f/g\to c\in(0,\infty)$ gives the same conclusion. Absolute convergence implies convergence. Dirichlet: if $F(R)=\int_a^Rf$ is bounded and $g$ decreases to zero, then $\int_a^\infty fg$ converges; Abel: if $\int_a^\infty f$ converges and $g$ is monotone bounded, then $\int_a^\infty fg$ converges.}
\EX{$\int_1^\infty\sin x/x\,dx$ 条件收敛：$|\int_1^R\sin x\,dx|\le2$ 且 $1/x\downarrow0$，Dirichlet 判别；但 $\int_1^\infty|\sin x|/x\,dx$ 发散（按周期分组）。$\int_{-1}^1dx/x$ 发散，Cauchy 主值为 $0$，主值有限不能替代收敛。}{$\int_1^\infty\sin x/x\,dx$ converges conditionally: $|\int_1^R\sin x\,dx|\le2$ and $1/x\downarrow0$, so Dirichlet applies; but $\int_1^\infty|\sin x|/x\,dx$ diverges by grouping over periods. $\int_{-1}^1dx/x$ diverges while its Cauchy principal value is $0$; a finite principal value does not imply convergence.}
\REMARK{几何应用（累加原则）：面积 $\int_a^b|f-g|$、绕 $x$ 轴体积 $\pi\int f^2$、弧长 $\int\sqrt{1+(f')^2}$、曲率 $|f''|/(1+(f')^2)^{3/2}$。数值积分：复合梯形误差 $O(h^2)$（$f\in C^2$），复合 Simpson 误差 $O(h^4)$（$f\in C^4$，$n$ 偶），后者对三次多项式精确；误差界要求光滑性，有尖点或奇异性时不能直接套用。}{Geometric applications (accumulation): area $\int_a^b|f-g|$, volume of revolution $\pi\int f^2$, arc length $\int\sqrt{1+(f')^2}$, curvature $|f''|/(1+(f')^2)^{3/2}$. Quadrature: composite trapezoidal has error $O(h^2)$ for $f\in C^2$, composite Simpson $O(h^4)$ for $f\in C^4$ with even $n$ (exact for cubics); the bounds require smoothness and fail at cusps or singularities.}
\ASK{微积分基本定理两个方向的条件是什么？}{What are the hypotheses of the two directions of the fundamental theorem?}{第一方向：$f$ 可积给 $F$ 连续，$f$ 在一点连续给 $F'=f$（对积分求导）。第二方向：$f$ 可积且有原函数 $F$ 时 $\int_a^bf=F(b)-F(a)$（用原函数求积分）。}{First direction: integrability gives continuity of $F$, and continuity of $f$ at a point gives $F'=f$ (differentiating an integral). Second: integrability plus an antiderivative gives $\int_a^bf=F(b)-F(a)$ (evaluating by an antiderivative).}

\section{数项级数 | Numerical series}

\subsection{部分和、Cauchy 准则与必要条件 | Partial sums, Cauchy criterion, necessary condition}
\BOTH{级数 $\sum a_n$ 由部分和 $s_N=\sum_{n=1}^N a_n$ 的极限定义；收敛即部分和收敛，极限称和。Cauchy 准则：收敛当且仅当对每个 $\varepsilon>0$ 存在 $N$，使 $m>n\ge N$ 时 $|\sum_{k=n+1}^m a_k|<\varepsilon$。必要条件：收敛必有 $a_n\to0$（$a_n=s_n-s_{n-1}$），反之不成立（调和级数）。}{A series is the limit of its partial sums $s_N=\sum_{n=1}^N a_n$; the limit is the sum. Cauchy criterion: convergence iff tails $|\sum_{k=n+1}^m a_k|$ can be made $<\varepsilon$ for $m>n\ge N$. Necessary condition: convergence forces $a_n\to0$ ($a_n=s_n-s_{n-1}$), not conversely (harmonic series).}
\EX{几何级数 $\sum r^n$：$r\ne1$ 时 $s_N=(1-r^N)/(1-r)$，故 $|r|<1$ 收敛且和为 $1/(1-r)$，$|r|\ge1$ 发散（通项不趋于零）。$p$-级数 $\sum1/n^p$ 在 $p>1$ 收敛、$p\le1$ 发散；$p=1$ 用二进分组（每组和至少 $1/2$），$p>1$ 用积分判别。}{For the geometric series $\sum r^n$: if $r\ne1$, $s_N=(1-r^N)/(1-r)$, so it converges with sum $1/(1-r)$ for $|r|<1$ and diverges for $|r|\ge1$ (terms do not tend to zero). The $p$-series $\sum1/n^p$ converges for $p>1$ and diverges for $p\le1$; $p=1$ by dyadic grouping, $p>1$ by the integral test.}

\subsection{正项级数的判别法 | Tests for positive-term series}
\BOTH{正项级数（$a_n\ge0$）部分和单调递增，由单调有界定理收敛当且仅当部分和有界，这是实数完备性在级数论中的第一次使用。}{For positive series ($a_n\ge0$) the partial sums are nondecreasing, so by monotone convergence the series converges iff the partial sums are bounded; this is the first use of completeness in series theory.}
\BOTH{\textbf{比较判别法.} 若存在 $M>0$ 使 $0\le a_n\le Mb_n$，则 $\sum b_n$ 收敛推 $\sum a_n$ 收敛，$\sum a_n$ 发散推 $\sum b_n$ 发散；极限形式 $a_n/b_n\to c$：$0<c<\infty$ 同敛散，$c=0$ 由 $\sum b_n$ 收敛推 $\sum a_n$ 收敛，$c=\infty$ 由 $\sum b_n$ 发散推 $\sum a_n$ 发散。}{\textbf{Comparison test.} If $0\le a_n\le Mb_n$ for some $M>0$, then convergence of $\sum b_n$ implies convergence of $\sum a_n$, and divergence of $\sum a_n$ implies divergence of $\sum b_n$; limit form $a_n/b_n\to c$: $0<c<\infty$ gives the same behaviour, $c=0$ transfers convergence, $c=\infty$ transfers divergence.}
\BOTH{\textbf{比值判别法（D'Alembert）.} $a_n>0$、$a_{n+1}/a_n\to L$：$L<1$ 时从某项起 $a_{n+1}\le qa_n$（$q\in(L,1)$），被几何级数控制，收敛；$L>1$ 时通项不趋于零，发散；$L=1$ 失效。\textbf{根值判别法（Cauchy）.} $\sqrt[n]{a_n}\to L$（或 $\limsup$）：结论同上，适用范围更广（不要求比值极限存在）。\textbf{积分判别法（Cauchy--Maclaurin）.} $a_n=f(n)$，$f$ 连续非负递减：级数与 $\int_1^\infty f$ 同敛散。}{\textbf{Ratio test (D'Alembert).} $a_n>0$, $a_{n+1}/a_n\to L$: if $L<1$, eventually $a_{n+1}\le qa_n$ for some $q\in(L,1)$, dominated by a geometric sequence, so it converges; if $L>1$ the terms do not tend to zero and it diverges; $L=1$ is inconclusive. \textbf{Root test (Cauchy).} $\sqrt[n]{a_n}\to L$ (or $\limsup$): same conclusions, wider scope (no ratio limit needed). \textbf{Integral test (Cauchy--Maclaurin).} $a_n=f(n)$ with $f$ continuous, nonnegative, decreasing: the series and $\int_1^\infty f$ converge together.}
\EX{$\sum1/n^2$ 收敛：$n\ge2$ 时 $1/n^2\le1/(n(n-1))=1/(n-1)-1/n$，部分和上界 $2$。$\sum n/2^n$：比值 $(n+1)/(2n)\to1/2<1$ 收敛。$\sum n!/n^n$：比值 $(n/(n+1))^n\to1/e<1$ 收敛。}{$\sum1/n^2$ converges: for $n\ge2$, $1/n^2\le1/(n(n-1))=1/(n-1)-1/n$, bounding partial sums by $2$. $\sum n/2^n$: ratio $(n+1)/(2n)\to1/2<1$, so it converges. $\sum n!/n^n$: ratio $(n/(n+1))^n\to1/e<1$, so it converges.}
\ASK{比值判别极限为 $1$ 说明什么？}{What does a ratio-test limit of one mean?}{它不能判定敛散性：$\sum1/n$ 与 $\sum1/n^2$ 的比值极限都是 $1$，前者发散后者收敛。应改用比较、积分或更精细的判别法，并明确说“此判别法失效”而非“发散”。}{It is inconclusive: $\sum1/n$ and $\sum1/n^2$ both have ratio limit $1$, one diverging and the other converging. Switch to comparison, the integral test, or a finer test, and say “the test is inconclusive” rather than “diverges”.}

\subsection{任意项级数：Leibniz、绝对收敛、Dirichlet 与 Abel | General series}
\BOTH{\textbf{Leibniz 判别法.} $a_n>0$ 单调递减趋于零，则 $\sum(-1)^{n-1}a_n$ 收敛，余项 $|s-s_N|\le a_{N+1}$ 且符号同首个遗漏项。证：部分和的偶子列递增有上界、奇子列递减有下界，两者之差 $a_{2n+1}\to0$。单调性不能去（把奇项取 $1/n$、偶项取 $1/n^2$ 则发散）。}{\textbf{Leibniz's test.} If $a_n>0$ decreases to zero, then $\sum(-1)^{n-1}a_n$ converges with error $|s-s_N|\le a_{N+1}$ of the sign of the first omitted term. Proof: the even partial sums increase and are bounded, the odd ones decrease and are bounded, and their difference $a_{2n+1}\to0$. Monotonicity cannot be dropped (taking $1/n$ on odd indices and $1/n^2$ on even ones gives divergence).}
\BOTH{\textbf{绝对收敛与条件收敛.} $\sum a_n$ 绝对收敛若 $\sum|a_n|$ 收敛，条件收敛若 $\sum a_n$ 收敛而 $\sum|a_n|$ 发散。绝对收敛推出收敛（三角不等式 $|\sum_{k=n+1}^m a_k|\le\sum|a_k|$ + Cauchy 准则），且 $|\sum a_n|\le\sum|a_n|$。}{\textbf{Absolute and conditional convergence.} $\sum a_n$ is absolutely convergent if $\sum|a_n|$ converges, conditionally convergent if $\sum a_n$ converges but $\sum|a_n|$ diverges. Absolute convergence implies convergence (triangle inequality plus Cauchy) and $|\sum a_n|\le\sum|a_n|$.}
\BOTH{\textbf{Dirichlet 判别法.} $\sum a_n$ 部分和有界且 $b_n$ 单调趋于零，则 $\sum a_nb_n$ 收敛。\textbf{Abel 判别法.} $\sum a_n$ 收敛且 $b_n$ 单调有界，则 $\sum a_nb_n$ 收敛。二者基于 Abel 分部求和 $\sum_{n=m}^N a_nb_n=A_Nb_N-A_{m-1}b_m+\sum A_n(b_n-b_{n+1})$；Leibniz 是 Dirichlet 的特例（$a_n=(-1)^{n-1}$，部分和在 $\{0,1\}$）。}{\textbf{Dirichlet's test.} If the partial sums of $\sum a_n$ are bounded and $b_n$ decreases to zero, then $\sum a_nb_n$ converges. \textbf{Abel's test.} If $\sum a_n$ converges and $b_n$ is monotone bounded, then $\sum a_nb_n$ converges. Both rest on Abel summation by parts; Leibniz is Dirichlet with $a_n=(-1)^{n-1}$.}
\EX{交错调和级数 $\sum(-1)^{n-1}/n$ 条件收敛，和为 $\ln2$；按“一正两负”重排得 $\frac12\ln2$。$\sum\sin n/n$ 由 Dirichlet 收敛（$|\sum_{k=1}^n\sin k|\le1/\sin(1/2)$ 有界、$1/n\downarrow0$），但 $\sum|\sin n|/n$ 发散，故条件收敛。}{The alternating harmonic series converges conditionally with sum $\ln2$; rearranged “one positive, two negative” it gives $\frac12\ln2$. $\sum\sin n/n$ converges by Dirichlet ($|\sum_{k=1}^n\sin k|$ is bounded, $1/n\downarrow0$), but $\sum|\sin n|/n$ diverges, so it is conditionally convergent.}
\BOTH{运算与重排：收敛级数可逐项相加、数乘、改有限项、加括号（新部分和是原部分和的子列）；但“加括号后收敛”不推出原级数收敛（$1-1+1-1+\cdots$）。绝对收敛级数任意重排和不变；条件收敛级数正部与负部都发散，Riemann 重排定理说可重排成任意指定和（含 $\pm\infty$）。Cauchy 乘积：两个绝对收敛级数的 Cauchy 乘积绝对收敛且和等于和的乘积；仅条件收敛时乘积可能发散。}{Operations and rearrangement: convergent series can be added termwise, scaled, changed in finitely many terms, and bracketed (new partial sums form a subsequence); but “convergent after bracketing” does not imply convergence ($1-1+1-1+\cdots$). An absolutely convergent series has rearrangement-invariant sum; a conditionally convergent series has divergent positive and negative parts, and Riemann's theorem says it can be rearranged to any prescribed sum (including $\pm\infty$). The Cauchy product of two absolutely convergent series is absolutely convergent with sum equal to the product of the sums; for merely conditional series the product may diverge.}
\ASK{为什么绝对收敛比条件收敛更稳定？}{Why is absolute convergence more robust than conditional convergence?}{绝对值级数收敛时尾项绝对值总和可任意小，重排、加括号、乘法的结论都不变；条件收敛依赖正负项相消，重排改变抵消顺序和结果。条件收敛级数的正部与负部都发散，所以“和”是两个无穷大之差。}{When the absolute series converges, the total tail magnitude is arbitrarily small, so rearrangement, bracketing and products behave predictably. Conditional convergence depends on cancellation, and rearrangement changes the order and result; both the positive and negative parts diverge, so the sum is a difference of two infinite quantities.}

\section{多元极限与连续 | Limits and continuity in \texorpdfstring{$\R^n$}{R\^{}n}}

\subsection{范数、点列收敛与紧致性 | Norms, convergence and compactness}
\BOTH{$\R^n$ 上定义内积 $\langle x,y\rangle=\sum x_iy_i$、范数 $\|x\|=\sqrt{\langle x,x\rangle}$ 与距离；Cauchy--Schwarz 不等式 $|\langle x,y\rangle|\le\|x\|\|y\|$ 给出三角不等式。所有范数等价（$\|x\|_\infty\le\|x\|_2\le\sqrt n\,\|x\|_\infty$），故开集、闭集、紧致、极限与连续不依赖范数选择，只在估计常数时需注意。点列收敛 $x_k\to x$ 当且仅当每个坐标收敛，这把 $\R^n$ 的问题化为 $n$ 个一维问题，并给出 $\R^n$ 的完备性（Cauchy 列逐坐标收敛）。}{On $\R^n$ define the inner product, norm and distance; Cauchy--Schwarz gives the triangle inequality. All norms are equivalent ($\|x\|_\infty\le\|x\|_2$ and $\|x\|_2\le\sqrt n\,\|x\|_\infty$), so open sets, closed sets, compactness, limits and continuity are independent of the norm, only constants in estimates depend on it. Convergence $x_k\to x$ iff each coordinate converges, reducing $\R^n$ to $n$ one-dimensional problems and giving completeness of $\R^n$.}
\BOTH{开集（每点有含于其中的球）、闭集（补集开；等价于含所有收敛点列的极限）、闭包 $\overline S=S\cup S'$、有界。紧致用开覆盖定义：每个开覆盖都有有限子覆盖。在 $\R^n$ 中三者等价：紧致 $\Leftrightarrow$ 闭且有界 $\Leftrightarrow$ 序列紧致（每个点列有收敛子列且极限在集合内）。反例：$(0,1)$ 的开覆盖 $(1/n,1)$ 无有限子覆盖；无界集跑向无穷、不闭集有跑出集合的收敛点列。}{\textbf{Open} sets (each point has a ball inside), \textbf{closed} sets (complement open; equivalently closed under limits of sequences), closure $\overline S=S\cup S'$, bounded sets. Compactness is defined by open covers. In $\R^n$ three things are equivalent: compact, closed and bounded, and sequentially compact (every sequence has a convergent subsequence whose limit lies in the set). Counterexamples: $(0,1)$ with the cover $(1/n,1)$ has no finite subcover; an unbounded set escapes to infinity, a non-closed set has a sequence converging outside.}
\EX{闭圆盘 $D=\{x^2+y^2\le1\}$ 闭且有界，故紧致；开圆盘有界不闭；双曲线 $\{xy=1\}$ 闭但无界（点列 $(1/n,n)$ 无收敛子列）。$f(x,y)=3x+4y$ 在 $D$ 上连续故取到最大值，Cauchy--Schwarz 给 $f\le5\sqrt{x^2+y^2}\le5$，在 $(3/5,4/5)$ 取等，故最大值为 $5$：紧致保证存在，不等式给出具体值。}{The closed disk $D=\{x^2+y^2\le1\}$ is closed and bounded, hence compact; the open disk is bounded but not closed; the hyperbola $\{xy=1\}$ is closed but unbounded (the sequence $(1/n,n)$ has no convergent subsequence). $f(x,y)=3x+4y$ is continuous on $D$ so attains a maximum; Cauchy--Schwarz gives $f\le5$ with equality at $(3/5,4/5)$, so the maximum is $5$: compactness gives existence, the inequality the value.}

\subsection{多元极限、路径依赖与累次极限 | Limits, path dependence and iterated limits}
\BOTH{设 $a$ 是 $D$ 的聚点。$\lim_{x\to a}f(x)=L$ 指对每个 $\varepsilon>0$ 存在 $\delta>0$，使 $x\in D$、$0<\|x-a\|<\delta$ 时 $|f(x)-L|<\varepsilon$；一个 $\delta$ 同时控制所有方向。Heine 判据：极限为 $L$ 当且仅当每个 $x_k\to a$（$x_k\ne a$）都有 $f(x_k)\to L$。否证极限找两条路径给不同值；证明极限用方向一致估计（范数夹逼或极坐标只依赖 $r$）。}{Let $a$ be a cluster point of $D$. $\lim_{x\to a}f(x)=L$ means for every $\varepsilon>0$ there is $\delta>0$ with $|f(x)-L|<\varepsilon$ whenever $0<\|x-a\|<\delta$; one $\delta$ controls every direction. Heine's criterion: the limit is $L$ iff every $x_k\to a$ with $x_k\ne a$ has $f(x_k)\to L$. Disprove by two paths with different limits; prove by a direction-independent estimate (norm squeeze, or polar coordinates depending only on $r$).}
\EX{$f(x,y)=xy/(x^2+y^2)$（$f(0,0)=0$）沿 $y=kx$ 等于 $k/(1+k^2)$，随 $k$ 变化，故原点极限不存在，尽管两个偏导都存在。更强：$g=x^2y/(x^4+y^2)$ 沿所有直线趋于 $0$，沿抛物线 $y=kx^2$ 却等于 $k/(1+k^2)$，故“所有直线方向一致”仍不够。正面：$|x^2y/(x^2+y^2)|\le|y|\le\sqrt{x^2+y^2}\to0$，故极限为 $0$。}{For $f(x,y)=xy/(x^2+y^2)$ with $f(0,0)=0$, along $y=kx$ the value is $k/(1+k^2)$, depending on $k$, so the limit fails even though both partials exist. Stronger: $g=x^2y/(x^4+y^2)$ tends to $0$ along every line but equals $k/(1+k^2)$ along $y=kx^2$, so agreement on all lines is still insufficient. Positively: $|x^2y/(x^2+y^2)|\le|y|\le\sqrt{x^2+y^2}\to0$, so the limit is $0$.}
\BOTH{累次极限 $\lim_{y\to y_0}\lim_{x\to x_0}f$ 先固定一个变量再令另一个趋近，与重极限量词顺序不同。四种情形都存在：重极限存在而累次极限不存在（$x\sin(1/y)+y\sin(1/x)$）；重极限存在、一个累次极限存在而另一个不存在（$x\sin(1/y)$）；两个累次极限相等而重极限不存在（$xy/(x^2+y^2)$）；两个累次极限不等（重极限必不存在）。唯一正面结论：若重极限为 $A$ 且对每个 $y$ 内层极限存在，则累次极限也是 $A$。}{The iterated limit first fixes one variable and lets the other tend; its quantifier order differs from the joint limit. All four cases occur: joint limit exists while iterated limits fail ($x\sin(1/y)+y\sin(1/x)$); joint limit exists with one iterated limit existing and the other not ($x\sin(1/y)$); both iterated limits equal yet the joint limit fails ($xy/(x^2+y^2)$); the two iterated limits differ (so the joint limit cannot exist). The only positive result: if the joint limit is $A$ and the inner limits exist, the iterated limit is also $A$.}

\subsection{连续、紧集三大定理与一致连续 | Continuity, the three theorems, uniform continuity}
\BOTH{$f$ 在 $a$ 连续指 $\lim_{x\to a}f(x)=f(a)$；和、积、商（分母非零）、复合保持连续；向量值函数连续当且仅当每个分量连续。分别连续不等于联合连续：$xy/(x^2+y^2)$ 对每个变量单独连续，但在原点不联合连续。紧集上的连续函数三大性质：紧集映成紧集、有界、取到最大最小值（证：取函数值趋于上确界的点列，抽收敛子列，闭性保极限点在集内，连续性给最值）；连通集的连续像连通，故紧连通集上取遍最值之间一切值。}{\textbf{Continuity} at $a$ means $\lim_{x\to a}f(x)=f(a)$; sums, products, quotients (nonzero denominator) and compositions preserve continuity; a vector-valued map is continuous iff each component is. Separate continuity does not imply joint continuity: $xy/(x^2+y^2)$ is separately continuous but not jointly continuous at the origin. On a compact set a continuous function has three properties: compact sets map to compact sets, boundedness, and attained extrema (take a sequence whose values approach the supremum, extract a convergent subsequence, closedness keeps the limit inside, continuity gives the value). A continuous image of a connected set is connected.}
\BOTH{\textbf{Cantor 定理.} 紧集上的连续函数必一致连续。证：否则有 $\varepsilon_0>0$ 及 $x_k,y_k$ 使 $\|x_k-y_k\|\to0$ 但 $\|f(x_k)-f(y_k)\|\ge\varepsilon_0$；抽子列 $y_{k_j}\to y_0$，则 $x_{k_j}\to y_0$，连续性矛盾。非紧集不保证：$x^2$ 在 $\R$ 上不一致连续。}{\textbf{Cantor's theorem.} A continuous function on a compact set is uniformly continuous. Proof: otherwise $\varepsilon_0>0$ and $x_k,y_k$ satisfy $\|x_k-y_k\|\to0$ but $\|f(x_k)-f(y_k)\|\ge\varepsilon_0$; a subsequence $y_{k_j}\to y_0$ forces $x_{k_j}\to y_0$, contradicting continuity. Noncompact sets do not guarantee it: $x^2$ on $\R$.}
\BOTH{多元没有 L'H\^opital 法则：它依赖全序、单调性和一元中值定理，而 $\R^n$ 无全序、中值定理只有沿线段的形式、趋近方向无穷多。替代工具是范数夹逼、极坐标、Taylor 展开与序列判据。}{There is no multivariable l'H\^opital rule: it relies on a total order, monotonicity and the one-variable mean value theorem, none of which transfers to $\R^n$ (no total order, the mean value theorem only along segments, infinitely many directions). Replacements are norm squeezing, polar coordinates, Taylor expansion and the sequential criterion.}
\ASK{证明或反证多元极限分别怎么做？}{How do you prove or disprove a multivariable limit?}{反证：找两条路径（直线、抛物线、坐标轴）给不同极限。证明：给方向一致估计 $|f(x)-L|\le C\|x-a\|^\alpha$，或极坐标后只依赖 $r$；序列判据统一两者。累次极限相等只是重极限存在的弱证据，不能作证明。}{Disprove by two paths (lines, parabolas, axes) with different limits; prove by a direction-independent estimate $|f(x)-L|\le C\|x-a\|^\alpha$, or by polar coordinates depending only on $r$; the sequential criterion unifies both. Equality of iterated limits is only weak evidence for a joint limit, not a proof.}

\section{多元微分 | Differentiation of several variables}

\subsection{偏导、方向导数与全微分 | Partial, directional and total derivatives}
\BOTH{偏导 $\partial f/\partial x_i$ 固定其余变量。方向导数 $D_uf(a)=\lim_{t\to0^+}[f(a+tu)-f(a)]/t$（$u$ 单位）。\textbf{全微分.} 存在线性映射 $A$（行向量 $\nabla f(a)$）使 $f(a+h)=f(a)+Ah+o(|h|)$，即 $f(a+h)=f(a)+\nabla f(a)\cdot h+o(|h|)$，全微分 $df=\sum f_{x_i}dx_i$。可微推出：连续、偏导存在且等于 $\nabla f$ 分量、$D_uf(a)=\nabla f(a)\cdot u$；反之不成立。}{\textbf{Partial} derivatives fix all but one variable; the \textbf{directional derivative} is a one-sided limit along a unit direction $u$. \textbf{Total differentiability} means a linear map $A=\nabla f(a)$ with $f(a+h)=f(a)+Ah+o(|h|)$, equivalently $df=\sum f_{x_i}dx_i$. Differentiability implies continuity, existence of partials equal to the components of $\nabla f$, and the formula $D_uf(a)=\nabla f(a)\cdot u$; the converses fail.}
\BOTH{可微的充分条件：偏导在某邻域存在且在 $a$ 连续。证：按坐标方向 telescoping，每段用一元 Lagrange 中值定理，再以偏导连续换 $\xi_i$ 为 $a$，误差 $o(|h|)$。条件充分不必要：$(x^2+y^2)\sin(1/\sqrt{x^2+y^2})$ 在原点可微但偏导不连续。方向导数存在也不能推连续：$x^3y/(x^6+y^2)$ 在原点所有方向导数为 $0$，却沿 $y=x^3$ 恒为 $1/2$。}{Sufficient condition for differentiability: the partials exist in a neighbourhood and are continuous at $a$. Proof: telescope the increment along coordinates, apply the one-variable mean value theorem per segment, and use continuity to replace $\xi_i$ by $a$ with error $o(|h|)$. The condition is sufficient but not necessary: $(x^2+y^2)\sin(1/\sqrt{x^2+y^2})$ is differentiable at the origin but its partials are discontinuous there. Directional derivatives do not even give continuity: $x^3y/(x^6+y^2)$ has all directional derivatives zero at the origin but is $1/2$ along $y=x^3$.}
\EX{$f=xy/\sqrt{x^2+y^2}$（$f(0,0)=0$）：两个偏导在原点为 $0$ 且函数连续，但可微要求余项比 $xy/(x^2+y^2)$ 趋于零，而它沿 $y=x$ 等于 $1/2$，故不可微。}{For $f=xy/\sqrt{x^2+y^2}$ with $f(0,0)=0$, both partials vanish and the function is continuous, but differentiability requires the remainder ratio $xy/(x^2+y^2)$ to tend to zero, which it fails along $y=x$ (value $1/2$), so $f$ is not differentiable.}
\ASK{偏导存在、方向导数存在、可微三者是什么关系？}{How are partial derivatives, directional derivatives and differentiability related?}{可微最强，推出偏导、任意方向导数公式和连续；偏导存在最弱，$xy/\sqrt{x^2+y^2}$ 在原点偏导存在、函数连续却不可微；方向导数存在也不能推连续，$x^3y/(x^6+y^2)$ 在原点所有方向导数为 $0$ 却沿 $y=x^3$ 恒为 $1/2$。}{Differentiability is strongest: it gives partials, the directional-derivative formula and continuity. Existence of partials is weakest: $xy/\sqrt{x^2+y^2}$ has both partials and is continuous at the origin but not differentiable. Directional derivatives do not even give continuity: $x^3y/(x^6+y^2)$ has all directional derivatives zero at the origin yet is discontinuous.}

\subsection{Jacobian、链式法则与隐函数/反函数定理 | Jacobians, the chain rule, implicit and inverse functions}
\FORMULA{链式法则与 Jacobian}{the chain rule and the Jacobian}
\BOTH{向量值映射 $F:\R^n\to\R^m$ 可微指存在 $m\times n$ 矩阵 $A$ 使 $F(a+h)=F(a)+Ah+o(|h|)$，$A$ 即 Jacobian $J_F=(\partial F_i/\partial x_j)$。链式法则 $J_{F\circ G}(x)=J_F(G(x))J_G(x)$（右到左，先内层再外层）；写成分量 $\partial z/\partial x_i=\sum_j(\partial f/\partial u_j)(\partial u_j/\partial x_i)$。$m=n$ 时 $|\det J_F(a)|$ 是体积伸缩因子。极坐标、球坐标的 Jacobian 分别为 $r$ 与 $\rho^2\sin\varphi$。}{A vector-valued map $F$ is differentiable if $F(a+h)=F(a)+Ah+o(|h|)$ with the $m\times n$ Jacobian matrix $A=J_F=(\partial F_i/\partial x_j)$. The chain rule is $J_{F\circ G}(x)=J_F(G(x))J_G(x)$ (right to left, inner map first); in coordinates $\partial z/\partial x_i=\sum_j(\partial f/\partial u_j)(\partial u_j/\partial x_i)$. When $m=n$, $|\det J_F(a)|$ is the local volume scaling. Polar and spherical Jacobians are $r$ and $\rho^2\sin\varphi$.}
\BOTH{\textbf{隐函数定理.} $F$ 在 $(x_0,y_0)$ 附近连续可偏导、$F(x_0,y_0)=0$、$F_y(x_0,y_0)\ne0$，则存在 $x_0$ 邻域与唯一 $C^1$ 函数 $y=f(x)$ 使 $F(x,f(x))=0$、$f(x_0)=y_0$、$f'=-F_x/F_y$。几何：$\nabla F$ 是等值线法向量，$F_y\ne0$ 表示等值线不竖直可写成图像。向量形式：$F=(F_1,\ldots,F_m):\R^{n+m}\to\R^m$，若 $F_y=(\partial F_i/\partial y_j)$ 可逆，则局部解出 $y=g(x)$，$J_g=-(F_y)^{-1}F_x$。\textbf{反函数定理}是其特例：$C^1$ 映射 $T$ 在 $x_0$ 处 $\det J_T\ne0$，则局部有 $C^1$ 逆，且 $J_S(y_0)=[J_T(x_0)]^{-1}$。}{\textbf{Implicit function theorem.} If $F$ has continuous partials near $(x_0,y_0)$, $F(x_0,y_0)=0$ and $F_y(x_0,y_0)\ne0$, then near $x_0$ there is a unique $C^1$ function $y=f(x)$ with $F(x,f(x))=0$, $f(x_0)=y_0$ and $f'=-F_x/F_y$. Geometrically $\nabla F$ is normal to the level curve and $F_y\ne0$ means it is not vertical, so it is a graph. Vector form: for $F:\R^{n+m}\to\R^m$ with $F_y$ invertible, locally $y=g(x)$ with $J_g=-(F_y)^{-1}F_x$. The \textbf{inverse function theorem} is the special case: a $C^1$ map $T$ with $\det J_T\ne0$ at $x_0$ has a local $C^1$ inverse with $J_S(y_0)=[J_T(x_0)]^{-1}$.}
\EX{曲线 $x^3+y^3-3xy=0$：$F_x=3x^2-3y$、$F_y=3y^2-3x$，在 $(3/2,3/2)$ 处 $y'=-1$；$F_y=0$ 的点（$(0,0)$、$(\sqrt[3]4,\sqrt[3]2)$）定理失效。$F=y^3-x^2=0$ 在原点 $F_y=0$，等值线 $y=x^{2/3}$ 不可微，无 $C^1$ 隐函数。}{For the curve $x^3+y^3-3xy=0$, $F_x=3x^2-3y$ and $F_y=3y^2-3x$; at $(3/2,3/2)$ we get $y'=-1$, while the theorem fails at the points with $F_y=0$. For $F=y^3-x^2=0$ at the origin $F_y=0$, the level curve $y=x^{2/3}$ is not differentiable, so no $C^1$ implicit function exists.}

\subsection{梯度、Hessian、极值与 Lagrange 乘子 | Gradient, Hessian, extrema and Lagrange multipliers}
\BOTH{若 $\nabla f(a)\ne0$，$D_uf(a)=\nabla f(a)\cdot u=|\nabla f(a)|\cos\theta$，故梯度是增长最快方向、最大方向导数为 $|\nabla f|$，且梯度是等值面的法向量。$f\in C^2$ 时混合偏导相等（Clairaut），Hessian $H=(f_{x_ix_j})$ 对称；二阶 Taylor：$f(a+h)=f(a)+\nabla f(a)\cdot h+\tfrac12h^{\mathsf T}H(a)h+o(|h|^2)$，特征值是主方向的二阶变化率。}{If $\nabla f(a)\ne0$, then $D_uf(a)=\nabla f(a)\cdot u=|\nabla f(a)|\cos\theta$; the gradient is the direction of steepest increase, the maximum directional derivative is $|\nabla f|$, and the gradient is normal to level surfaces. For $f\in C^2$ the mixed partials agree (Clairaut) and the Hessian is symmetric; the second-order Taylor formula is as displayed, and the eigenvalues are the second-order rates along principal directions.}
\BOTH{\textbf{无条件极值.} 必要条件（Fermat）：内点局部极值且可微则 $\nabla f=0$（驻点）。充分条件（Hessian 判别）：$f\in C^2$、$\nabla f(a)=0$ 时，$H(a)$ 正定给严格局部极小、负定给严格局部极大、不定给鞍点；二元情形 $f_{xx}f_{yy}-f_{xy}^2>0$ 且 $f_{xx}>0$（$<0$）为极小（极大），行列式 $<0$ 为鞍点，$=0$ 失效。}{\textbf{Unconstrained extrema.} Necessary (Fermat): an interior local extremum of a differentiable function is stationary, $\nabla f=0$. Sufficient (Hessian test): for $f\in C^2$ with $\nabla f(a)=0$, positive definiteness of $H(a)$ gives a strict local minimum, negative definiteness a strict local maximum, indefiniteness a saddle. In two variables $f_{xx}f_{yy}-f_{xy}^2>0$ with $f_{xx}>0$ ($<0$) gives a minimum (maximum); a negative determinant a saddle; a zero determinant is inconclusive.}
\BOTH{\textbf{条件极值与 Lagrange 乘子.} 在约束 $g_k(x)=0$（$k=1,\ldots,m<n$）下极小化 $f$。若 $f,g_k\in C^1$、$x_0$ 是条件极值点、约束 Jacobian 满秩，则存在乘子 $\lambda_k$ 使 $\nabla f(x_0)=\sum_k\lambda_k\nabla g_k(x_0)$。几何：$\nabla f$ 落在约束法空间 $\mathrm{span}\{\nabla g_k\}$ 内，否则沿可行集还有改进方向。充分条件在切空间 $T=\{h:\nabla g_k\cdot h=0\}$ 上用 $H_xL$ 的定号，$L=f-\sum\lambda_kg_k$；二阶判别只在切子空间上做。$\lambda_k=\partial f^*/\partial c_k$ 是最优值对约束水平的敏感度。}{\textbf{Constrained extrema and Lagrange multipliers.} Minimize $f$ subject to $g_k(x)=0$. If $f,g_k\in C^1$, $x_0$ is a constrained extremum and the constraint Jacobian has full rank, then $\nabla f(x_0)=\sum_k\lambda_k\nabla g_k(x_0)$ for some multipliers; geometrically $\nabla f$ lies in the normal space spanned by the $\nabla g_k$, otherwise the objective could still improve along the feasible set. The sufficient condition uses the definiteness of $H_xL$ on the tangent space $T$, with $L=f-\sum\lambda_kg_k$. Also $\lambda_k=\partial f^*/\partial c_k$ is the sensitivity of the optimal value to the constraint level.}
\EX{$f(x,y)=x^2+xy+y^2$：梯度 $(2x+y,x+2y)$ 只在原点为零，$H=\begin{pmatrix}2&1\\1&2\end{pmatrix}$ 正定（配平方 $f=(x+y/2)^2+3y^2/4$），故原点为（全局）严格极小。求 $V=8xyz$ 在 $x^2+y^2+z^2=1$、$x,y,z>0$ 的最大：$L=8xyz-\lambda(x^2+y^2+z^2-1)$ 得 $x=y=z=1/\sqrt3$，$V=8\sqrt3/9$。}{For $f(x,y)=x^2+xy+y^2$, the gradient vanishes only at the origin and $H$ is positive definite (completing squares), so the origin is a strict (global) minimum. Maximizing $V=8xyz$ on $x^2+y^2+z^2=1$ with $x,y,z>0$: from $L=8xyz-\lambda(x^2+y^2+z^2-1)$ we get $x=y=z=1/\sqrt3$ and $V=8\sqrt3/9$.}
\ASK{为什么约束需要 Jacobian 满秩？}{Why must the constraint Jacobian have full rank?}{满秩保证 $\nabla g_k$ 线性无关，可行集在极值点附近是光滑的 $(n-m)$ 维流形，切空间与法空间有意义；否则可行集退化，Lagrange 条件失效。例：在 $g=x^2=0$ 下最小化 $f=x$，解为 $(0,0)$，但 $\nabla g=(0,0)$，不存在满足 $\nabla f=\lambda\nabla g$ 的乘子。}{Full rank makes the normals independent, so the feasible set is a smooth $(n-m)$-manifold near the extremum and tangent/normal spaces are defined; otherwise the set degenerates and the Lagrange condition fails. Example: minimize $f=x$ subject to $g=x^2=0$; the solution is $(0,0)$, but $\nabla g=(0,0)$ and no multiplier satisfies $\nabla f=\lambda\nabla g$.}

\section{多元积分 | Integration of several variables}

\subsection{重积分、Fubini 与变量替换 | Multiple integrals, Fubini and change of variables}
\BOTH{有界区域上对小块作 Riemann 和 $\sum f(p_i)\operatorname{area}(D_i)$，极限与分割取点无关时称可积，记 $\iint_Df\,dA$。有界闭区域上连续函数可积，$f\equiv1$ 给面积，面密度给质量。性质：线性、区域可加、保序、$|\iint f|\le\iint|f|$、积分中值定理。}{\textbf{Multiple integrals.} For a bounded region, Riemann sums over small pieces define the integral if the limit is independent of partition and sample points; $f\equiv1$ gives area and a surface density gives mass. Properties: linearity, additivity, monotonicity, $|\iint f|\le\iint|f|$, and the mean value theorem.}
\BOTH{\textbf{Fubini 定理.} 闭矩形上若 $f$ 可积且每个 $f(x,\cdot)$ 可积，则 $\iint f\,dA=\int_a^b(\int_c^df(x,y)dy)dx$，交换次序同理；连续函数满足条件。条件不能省：$(x^2-y^2)/(x^2+y^2)^2$ 在 $(0,1)^2$ 上两个累次积分分别为 $-\pi/4$ 与 $+\pi/4$，因 $\iint|f|=\infty$；非负函数用 Tonelli 恒可交换。}{Fubini's theorem: on a rectangle, if $f$ is integrable and each $f(x,\cdot)$ is integrable, the double integral equals the iterated integral in either order; continuous functions satisfy all hypotheses. The hypothesis is essential: $(x^2-y^2)/(x^2+y^2)^2$ on $(0,1)^2$ has iterated integrals $-\pi/4$ and $+\pi/4$, since $\iint|f|=\infty$; nonnegative functions always permit Fubini via Tonelli.}
\FORMULA{变量替换与 Jacobian}{change of variables and the Jacobian}
\BOTH{设 $T:D\to\Omega$ 是 $C^1$ 一一映射，$f$ 连续，则
\[
\iint_\Omega f(x,y)\,dx\,dy=\iint_D f(T(u,v))\,|\det J_T(u,v)|\,du\,dv .
\]
条件：一一（否则重复计数）、$C^1$、Jacobian 除零面积集外非零。极坐标 $|\det J|=r$，球坐标 $|\det J|=\rho^2\sin\varphi$（$\varphi$ 是与 $z$ 轴夹角）。}{If $T:D\to\Omega$ is a one-to-one $C^1$ map and $f$ is continuous, the displayed change-of-variables formula holds. Hypotheses: injectivity (else double counting), $C^1$, nonzero Jacobian off a null set. Polar Jacobian is $r$; spherical is $\rho^2\sin\varphi$ (with $\varphi$ the angle from the $z$-axis).}
\EX{极坐标算 $I=\iint_{x^2+y^2\le R^2}e^{-(x^2+y^2)}dA=\int_0^{2\pi}\int_0^R e^{-r^2}r\,dr\,d\theta=\pi(1-e^{-R^2})$；令 $R\to\infty$ 得 Gauss 积分 $\int_{-\infty}^{\infty}e^{-x^2}dx=\sqrt\pi$。反常重积分：$\iint_{r\ge1}r^{-p}dA=2\pi\int_1^\infty r^{1-p}dr$ 当且仅当 $p>2$ 收敛。}{In polar coordinates $I=\int_0^{2\pi}\int_0^R e^{-r^2}r\,dr\,d\theta=\pi(1-e^{-R^2})$; letting $R\to\infty$ gives the Gaussian integral $\int_{-\infty}^{\infty}e^{-x^2}dx=\sqrt\pi$. For improper double integrals, $\iint_{r\ge1}r^{-p}dA=2\pi\int_1^\infty r^{1-p}dr$ converges exactly for $p>2$.}

\subsection{两类曲线积分与曲面积分 | Line and surface integrals of two kinds}
\BOTH{\textbf{第一型（弧长/面积）} 与方向无关：$\int_Cf\,ds=\int_a^bf(r(t))|r'(t)|dt$；$\iint_Sf\,dS=\iint_Df(r(u,v))|r_u\times r_v|du\,dv$。\textbf{第二型（功/通量）} 依赖定向：$\int_CF\cdot dr=\int_a^bF(r(t))\cdot r'(t)dt$；$\iint_SF\cdot n\,dS=\iint_DF(r(u,v))\cdot(r_u\times r_v)du\,dv$。联系：$dr=\tau\,ds$、$n\,dS=r_u\times r_v\,du\,dv$，故第二型是第一型对切向/法向分量的积分，反向变号。}{The \textbf{first kind} (arc length / area) is orientation-free: $\int_Cf\,ds=\int_a^bf(r(t))|r'(t)|dt$ and $\iint_Sf\,dS=\iint_Df(r(u,v))|r_u\times r_v|du\,dv$. The \textbf{second kind} (work / flux) depends on orientation: $\int_CF\cdot dr=\int_a^bF(r(t))\cdot r'(t)dt$ and $\iint_SF\cdot n\,dS=\iint_DF(r(u,v))\cdot(r_u\times r_v)du\,dv$. The relation is $dr=\tau\,ds$ and $n\,dS=r_u\times r_v\,du\,dv$, so the second kind is the first kind applied to the tangential or normal component, changing sign under reversal.}
\EX{沿 $r(t)=(t,t)$、$0\le t\le1$：$\int_Cx\,ds=\sqrt2/2$ 而 $\int_Cx\,dx=1/2$，两类累加不同量。沿单位圆逆时针 $F=(-y,x)$ 的环流为 $2\pi$（顺时针 $-2\pi$）；$F=(x,y)$ 从 $(0,0)$ 到 $(1,1)$ 沿直线与折线都得 $1$，因 $F=\nabla(\tfrac12(x^2+y^2))$，积分只依赖端点。}{Along $r(t)=(t,t)$ for $0\le t\le1$, $\int_Cx\,ds=\sqrt2/2$ but $\int_Cx\,dx=1/2$: the two kinds accumulate different quantities. The circulation of $F=(-y,x)$ around the counterclockwise unit circle is $2\pi$ (clockwise $-2\pi$); for $F=(x,y)$ from $(0,0)$ to $(1,1)$ both the segment and a broken line give $1$, since $F=\nabla(\tfrac12(x^2+y^2))$ and the integral depends only on the endpoints.}

\subsection{Green、Gauss、Stokes 公式与场论 | Green, Gauss, Stokes and field theory}
\BOTH{\textbf{Green 公式.} $D$ 由分片光滑 Jordan 曲线围成、$\partial D$ 正向（区域在左）、$P,Q\in C^1(\overline D)$，则 $\oint_{\partial D}P\,dx+Q\,dy=\iint_D(Q_x-P_y)dA$。\textbf{Gauss（散度）公式.} 有界 $\Omega\subset\R^3$、边界分片光滑闭曲面外法向、$F\in C^1(\overline\Omega)$，则 $\iint_{\partial\Omega}F\cdot n\,dS=\iiint_\Omega\nabla\cdot F\,dV$。\textbf{Stokes 公式.} 光滑定向曲面 $\Sigma$、边界取诱导定向、$F\in C^1$，则 $\oint_{\partial\Sigma}F\cdot dr=\iint_\Sigma(\nabla\times F)\cdot n\,dS$。}{Green's theorem: for $D$ bounded by piecewise smooth Jordan curves with positive boundary and $P,Q\in C^1(\overline D)$, $\oint_{\partial D}P\,dx+Q\,dy=\iint_D(Q_x-P_y)dA$. Gauss's (divergence) theorem: for bounded $\Omega\subset\R^3$ with piecewise smooth closed boundary, outward normal and $F\in C^1(\overline\Omega)$, $\iint_{\partial\Omega}F\cdot n\,dS=\iiint_\Omega\nabla\cdot F\,dV$. Stokes's theorem: for a smooth oriented surface $\Sigma$ with induced boundary orientation and $F\in C^1$, $\oint_{\partial\Sigma}F\cdot dr=\iint_\Sigma(\nabla\times F)\cdot n\,dS$.}
\BOTH{定向不可省：反向差负号；曲面须可定向（Möbius 带不可）。Green 是 Stokes 的平面特例；三者共同思想是边界上的积分等于内部导数的积分。条件反例：$F=r/|r|^3$ 在原点无定义，散度处处为零却通过含原点的闭曲面通量为 $4\pi$，Gauss 失效（$C^1$ 条件破坏）。}{Orientation is essential (reversal changes the sign) and surfaces must be orientable (not the Möbius strip). Green is the planar case of Stokes; all three express a boundary integral as an integral of an interior derivative. Counterexample: $F=r/|r|^3$ is undefined at the origin with zero divergence away from it, yet has flux $4\pi$ through any surface enclosing the origin, so Gauss fails (the $C^1$ hypothesis is broken).}
\BOTH{\textbf{场论.} 散度 $\nabla\cdot F=P_x+Q_y+R_z$ 是通量体密度（源正汇负，$\nabla\cdot F=0$ 无源/不可压缩）；旋度 $\nabla\times F$ 是环量面密度（$\nabla\times F=0$ 无旋/保守）。恒等式 $\nabla\cdot(\nabla\times F)=0$、$\nabla\times(\nabla f)=0$。在单连通区域上 $F\in C^1$ 时，路径无关、存在势 $F=\nabla U$、$\nabla\times F=0$ 三者等价；平面即 $Q_x=P_y$。反例 $F=(-y/(x^2+y^2),x/(x^2+y^2))$ 在去心平面上旋度为零，但绕单位圆积分为 $2\pi$（区域有洞）。}{\textbf{Field theory.} Divergence $\nabla\cdot F$ is the volume density of flux (sources positive, sinks negative; $\nabla\cdot F=0$ incompressible); curl $\nabla\times F$ is the surface density of circulation ($\nabla\times F=0$ irrotational/conservative). The identities $\nabla\cdot(\nabla\times F)=0$ and $\nabla\times(\nabla f)=0$ hold. On a simply connected region with $F\in C^1$, path independence, existence of a potential $F=\nabla U$, and $\nabla\times F=0$ are equivalent ($Q_x=P_y$ in the plane). Counterexample: $F=(-y/(x^2+y^2),x/(x^2+y^2))$ has zero curl on the punctured plane but circulation $2\pi$ around the unit circle (the region has a hole).}
\EX{Green：$D$ 为三角形 $(0,0)\to(2,0)\to(2,1)$，$P=x^2-y$、$Q=x+y^2$，$Q_x-P_y=2$ 且面积 $1$，故环流 $2$；逐段计算同样得 $\frac83+\frac73-3=2$。Gauss：$F=(xy,yz,zx)$ 在 $[0,1]^3$ 上散度为 $y+z+x$，其积分为 $3/2$，直接算三个 $x=1,y=1,z=1$ 面各贡献 $1/2$ 一致。}{Green: for the triangle $(0,0)\to(2,0)\to(2,1)$ with $P=x^2-y$, $Q=x+y^2$, we have $Q_x-P_y=2$ and area $1$, so the circulation is $2$; direct edge-by-edge computation also gives $\frac83+\frac73-3=2$. Gauss: for $F=(xy,yz,zx)$ on $[0,1]^3$ the divergence is $y+z+x$ with integral $3/2$, matching the direct flux through the faces $x=1,y=1,z=1$ of $1/2$ each.}
\REMARK{守恒律：对任意控制体 $\Omega$，质量变化率等于负的边界通量，由 Gauss 公式得连续性方程 $\rho_t+\nabla\cdot(\rho v)=0$；热传导 $\rho c\,T_t=\nabla\cdot(k\nabla T)$ 同理。三个公式统一为一般 Stokes 公式 $\int_{\partial\Omega}\omega=\int_\Omega d\omega$。}{Conservation laws: the rate of change of mass in a control volume equals minus the boundary flux, which by Gauss gives the continuity equation $\rho_t+\nabla\cdot(\rho v)=0$; heat conduction $\rho cT_t=\nabla\cdot(k\nabla T)$ is analogous. The three theorems unify as the general Stokes formula $\int_{\partial\Omega}\omega=\int_\Omega d\omega$.}
\ASK{Gauss 公式为什么要求闭曲面与外法向？}{Why does Gauss require a closed surface and an outward normal?}{散度定理算净流出量：只有闭曲面区分内外，外法向使 $F\cdot n>0$ 表示流出。开曲面通量只是总通量的一部分；取内法向每项差负号。}{The divergence theorem computes net outflow: only a closed surface separates inside from outside, and the outward normal makes $F\cdot n>0$ mean outflow. An open surface gives only part of the total flux; an inward normal flips every sign.}

\section{含参积分 | Integrals depending on a parameter}

\subsection{连续性、可导性与可积性 | Continuity, differentiability and integrability}

\BOTH{设 $f$ 在矩形 $R=[a,b]\times[c,d]$ 上连续，$I(y)=\int_a^b f(x,y)\,dx$ 称为\textbf{含参积分}，其中 $x$ 是积分变量、$y$ 是参数。\textbf{连续性.} $f$ 在紧矩形 $R$ 上一致连续：对每个 $\varepsilon>0$ 存在 $\delta>0$，使 $|y-y_0|<\delta$ 时对一切 $x\in[a,b]$ 有 $|f(x,y)-f(x,y_0)|<\varepsilon$，于是 $|I(y)-I(y_0)|\le\varepsilon(b-a)$，故 $I$ 在 $[c,d]$ 上连续。}{\textbf{Continuity.} Let $f$ be continuous on the rectangle $R=[a,b]\times[c,d]$; $I(y)=\int_a^b f(x,y)\,dx$ is an integral depending on the parameter $y$. Since $f$ is uniformly continuous on the compact rectangle, for every $\varepsilon>0$ there is $\delta>0$ such that $|y-y_0|<\delta$ implies $|f(x,y)-f(x,y_0)|<\varepsilon$ for all $x\in[a,b]$; hence $|I(y)-I(y_0)|\le\varepsilon(b-a)$ and $I$ is continuous on $[c,d]$.}

\FORMULA{积分号下求导与变限公式}{differentiation under the integral sign and variable limits}
\[
\frac{d}{dy}\int_{\alpha(y)}^{\beta(y)}f(x,y)\,dx
=f(\beta(y),y)\beta'(y)-f(\alpha(y),y)\alpha'(y)+\int_{\alpha(y)}^{\beta(y)}f_y(x,y)\,dx .
\]

\BOTH{若 $f$ 与 $f_y$ 在 $R$ 上连续，则 $I$ 在 $[c,d]$ 可导且 $I'(y)=\int_a^b f_y(x,y)\,dx$（Leibniz 法则）。证：差商等于 $\int_a^b f_y(x,y+\theta h)\,dx$（对固定的 $x$ 用 Lagrange 中值定理，$0<\theta<1$ 依赖 $x$ 与 $h$），再由 $f_y$ 在紧矩形上的一致连续性令 $h\to0$。上下限也依赖参数时，再加上两个边界项，即得上式。}{\textbf{Leibniz's rule.} If $f$ and $f_y$ are continuous on $R$, then $I$ is differentiable on $[c,d]$ with $I'(y)=\int_a^b f_y(x,y)\,dx$. Proof: the difference quotient equals $\int_a^b f_y(x,y+\theta h)\,dx$ for some $\theta\in(0,1)$ (Lagrange in $x$); uniform continuity of $f_y$ on the compact rectangle lets $h\to0$. When the limits depend on $y$, the two boundary terms are added, giving the displayed formula.}

\BOTH{\textbf{积分号下积分（交换次序）.} 若 $f$ 在 $R$ 上连续，则两个累次积分相等：$\int_c^d dy\int_a^bf(x,y)\,dx=\int_a^b dx\int_c^df(x,y)\,dy$。这是 Fubini 定理的连续函数情形；对含参反常积分，下面的一致收敛取代矩形上的连续性。}{\textbf{Integration under the integral sign.} If $f$ is continuous on $R$, the two iterated integrals agree: $\int_c^d dy\int_a^bf(x,y)\,dx=\int_a^b dx\int_c^df(x,y)\,dy$. This is the continuous case of Fubini's theorem; for improper parametric integrals, uniform convergence below replaces continuity on the rectangle.}

\EX{$I(y)=\int_0^1x^y\,dx$（$y>-1$）直接算得 $1/(y+1)$；由 Leibniz 法则 $I'(y)=\int_0^1x^y\log x\,dx=-1/(y+1)^2$（$y>0$ 时 $x^y\log x$ 在 $x=0$ 处补充为 $0$ 后连续）。但 $\int_0^1\frac{x^2-y^2}{(x^2+y^2)^2}\,dx$ 在 $y=0$ 处有瑕点、被积函数不连续，不能对每个 $y$ 直接求导，说明紧矩形上的连续性条件要紧。}{For $I(y)=\int_0^1x^y\,dx$ with $y>-1$, direct computation gives $1/(y+1)$; Leibniz's rule gives $I'(y)=\int_0^1x^y\log x\,dx=-1/(y+1)^2$ (for $y>0$, with $x^y\log x$ defined as $0$ at $x=0$). However $\int_0^1\frac{x^2-y^2}{(x^2+y^2)^2}\,dx$ has a singular point at $y=0$, so one cannot differentiate pointwise: the continuity hypothesis on a rectangle is essential.}

\ASK{什么时候可以在积分号下求导、交换积分次序？}{When may we differentiate under the integral sign or interchange integrals?}{求导需 $f$ 与 $f_y$ 在含参数的紧矩形上连续（含参反常积分则需 $\int f_y$ 一致收敛），从而差商被积函数一致趋于 $f_y$；交换次序需 $f$ 在矩形上连续（非负可测函数用 Tonelli）。只知每个 $y$ 的积分存在是不够的。}{Differentiation needs $f$ and $f_y$ continuous on a compact parameter rectangle (or uniform convergence of $\int f_y$ in the improper case), so the difference quotients converge uniformly to $f_y$; interchanging needs continuity of $f$ on the rectangle (Tonelli for nonnegative measurable functions). Existence of the integral for each $y$ alone is insufficient.}

\subsection{含参反常积分与一致收敛 | Improper parametric integrals and uniform convergence}

\BOTH{设 $f$ 在 $[a,\infty)\times[c,d]$ 上连续，$I(y)=\int_a^\infty f(x,y)\,dx$。称 $I$ 在 $[c,d]$ 上\textbf{一致收敛}，若对每个 $\varepsilon>0$ 存在 $R_0$（只依赖 $\varepsilon$），使 $R>R_0$ 时对每个 $y\in[c,d]$ 都有 $|\int_R^\infty f(x,y)\,dx|<\varepsilon$；等价于 $\sup_{y\in[c,d]}|\int_R^\infty f(x,y)\,dx|\to0$（$R\to\infty$）。}{\textbf{Uniform convergence.} Let $f$ be continuous on $[a,\infty)\times[c,d]$ and $I(y)=\int_a^\infty f(x,y)\,dx$. We say $I$ converges uniformly on $[c,d]$ if for every $\varepsilon>0$ there is $R_0$ (depending only on $\varepsilon$) such that $R>R_0$ implies $|\int_R^\infty f(x,y)\,dx|<\varepsilon$ for every $y\in[c,d]$; equivalently the supremum over $y$ of the tails tends to $0$.}

\BOTH{\textbf{Weierstrass（M）判别法.} 若 $|f(x,y)|\le g(x)$ 且 $\int_a^\infty g(x)\,dx$ 收敛，则含参积分在 $[c,d]$ 上一致收敛。Dirichlet 与 Abel 判别法也有含参版本，关键是部分和（或单调因子）的估计与 $y$ 无关。一致收敛保留连续性、积分号下积分与求导（求导还需 $\int f_y$ 一致收敛）；通常只须在每个闭子区间 $[c',d']\subset(c,d)$ 上验证，结论也是局部的。}{\textbf{Weierstrass (M) test.} If $|f(x,y)|\le g(x)$ and $\int_a^\infty g(x)\,dx$ converges, then the parametric integral converges uniformly on $[c,d]$. Dirichlet's and Abel's tests also have parametric versions; the key is that the bounds on partial sums (or on the monotone factor) do not depend on $y$. Uniform convergence preserves continuity, integration under the integral sign and differentiation (the last also needs uniform convergence of $\int f_y$); it is enough to check it on each closed subinterval $[c',d']\subset(c,d)$, and the conclusions are local.}

\EX{$\int_0^\infty e^{-xy}\,dx=1/y$（$y>0$）：在 $[c,d]\subset(0,\infty)$ 上由 $e^{-xy}\le e^{-cx}$ 控制，而 $\int_0^\infty e^{-cx}dx$ 收敛，故一致收敛；对 $y$ 求导得 $\int_0^\infty xe^{-xy}dx=1/y^2$，与直接计算一致。$y\to0^+$ 时积分发散，故结论只在 $(0,\infty)$ 的局部成立。标准例子是 $\Gamma(s)=\int_0^\infty x^{s-1}e^{-x}\,dx$：它在 $s>0$ 上连续可导，且 $\Gamma(s+1)=s\Gamma(s)$。}{$\int_0^\infty e^{-xy}\,dx=1/y$ for $y>0$: on $[c,d]\subset(0,\infty)$ it is dominated by $e^{-cx}$ with convergent integral $\int_0^\infty e^{-cx}dx$, hence uniformly convergent; differentiating gives $\int_0^\infty xe^{-xy}dx=1/y^2$, matching direct computation. The integral diverges as $y\to0^+$, so the conclusion is local to $(0,\infty)$. The standard example is $\Gamma(s)=\int_0^\infty x^{s-1}e^{-x}\,dx$, continuous and differentiable for $s>0$ with $\Gamma(s+1)=s\Gamma(s)$.}

\ASK{一致收敛与逐点收敛的差别在哪里？}{What is the difference between uniform and pointwise convergence?}{逐点收敛只保证每个参数点处尾部误差最终变小，$R_0$ 可以依赖 $y$；一致收敛要求 $R_0$ 与 $y$ 无关，因此极限、积分、求导才能与积分号交换。反例：$y^n$ 在 $[0,1)$ 上逐点趋于 $0$ 但不一致；含参积分的尾部估计也必须对 $y$ 一致。}{Pointwise convergence only makes the tail error eventually small at each parameter value, with $R_0$ allowed to depend on $y$; uniform convergence requires $R_0$ independent of $y$, which is what lets limits, integrals and derivatives pass under the integral sign. A counterexample is $y^n$ on $[0,1)$, which converges pointwise to $0$ but not uniformly; tail estimates for parametric integrals must also be uniform in $y$.}

\section{Fourier 级数 | Fourier series}

\subsection{正交系、系数、Bessel 与收敛定理 | Orthogonality, coefficients, Bessel and convergence}
\FORMULA{Fourier 系数与级数}{Fourier coefficients and series}
\BOTH{在 $[-\pi,\pi]$ 上取内积 $\langle f,g\rangle=\int_{-\pi}^{\pi}fg$，不同频率正余弦两两正交。若 $f\sim\frac{a_0}{2}+\sum(a_n\cos nx+b_n\sin nx)$，正交投影得
\[
a_n=\frac1\pi\int_{-\pi}^{\pi}f\cos(nx)\,dx,\quad
b_n=\frac1\pi\int_{-\pi}^{\pi}f\sin(nx)\,dx,\quad
a_0=\frac1\pi\int_{-\pi}^{\pi}f\,dx .
\]
周期 $2L$ 时频率 $n\pi/L$、区间 $[-L,L]$、前因子 $1/L$。}{On $[-\pi,\pi]$ with inner product $\langle f,g\rangle=\int_{-\pi}^{\pi}fg$, sines and cosines of distinct frequencies are orthogonal. Orthogonal projection gives the displayed coefficients. For period $2L$ use frequency $n\pi/L$, interval $[-L,L]$ and prefactor $1/L$.}
\BOTH{\textbf{Bessel 不等式.} 对任意 $N$，$\frac{a_0^2}{2}+\sum_{n=1}^{N}(a_n^2+b_n^2)\le\frac1\pi\int_{-\pi}^{\pi}f^2$。$S_N$ 在所有 $N$ 次三角多项式中使均方误差最小（正交投影/最佳逼近）。\textbf{Dirichlet--Jordan 收敛定理.} $f$ 以 $2\pi$ 为周期、可积；若 $f$ 在 $x$ 邻域分段单调或满足 Hölder 条件，则 Fourier 级数在 $x$ 收敛到 $\frac{f(x+)+f(x-)}{2}$（连续点为 $f(x)$，跳跃点为左右极限平均）。仅连续不足以保证逐点收敛（du Bois-Reymond 反例）。}{Bessel's inequality: for every $N$, $\frac{a_0^2}{2}+\sum_{n=1}^{N}(a_n^2+b_n^2)\le\frac1\pi\int_{-\pi}^{\pi}f^2$. The partial sum $S_N$ minimizes the mean-square error among trigonometric polynomials of degree $N$ (orthogonal projection). Dirichlet--Jordan theorem: for $2\pi$-periodic integrable $f$ that is piecewise monotone near $x$ or Hölder there, the Fourier series converges at $x$ to $\frac{f(x+)+f(x-)}{2}$ (to $f(x)$ at continuity points, the average of one-sided limits at jumps). Continuity alone does not ensure pointwise convergence (du Bois-Reymond).}
\BOTH{半区间 $[0,L]$ 上偶延拓给余弦级数（有常数项），奇延拓给正弦级数（无常数项）。函数越光滑，系数衰减越快（分段 $C^1$ 时 $b_n=O(1/n)$）。}{On $[0,L]$ an even extension gives a cosine series (with constant term), an odd extension a sine series (without). Smoother functions have faster coefficient decay (piecewise $C^1$ gives $b_n=O(1/n)$).}

\subsection{Parseval 等式与 Gibbs 现象 | Parseval's identity and the Gibbs phenomenon}
\FORMULA{Parseval 等式}{Parseval's identity}
\[
\frac{a_0^2}{2}+\sum_{n=1}^{\infty}(a_n^2+b_n^2)=\frac1\pi\int_{-\pi}^{\pi}f(x)^2\,dx .
\]
\BOTH{Parseval 等式说均方“能量”等于各频率分量能量之和，是 Bessel 不等式的极限情形，也是函数空间中的勾股定理。}{\textbf{Parseval's identity} says the mean-square energy equals the sum of the energies of the frequency components; it is the limiting case of Bessel and Pythagoras' theorem in function space.}
\EX{方波 $f=-1$（$-\pi<x<0$）、$f=1$（$0<x<\pi$），奇函数故 $a_n=0$，$b_n=4/(n\pi)$（$n$ 奇），得 $f\sim\frac4\pi\sum_{k\ge0}\frac{\sin((2k+1)x)}{2k+1}$；在 $x=\pi/2$ 得 Leibniz 公式 $\pi/4=1-1/3+1/5-\cdots$。Parseval 给出 $\sum_{k\ge0}1/(2k+1)^2=\pi^2/8$，进而 $\sum1/n^2=\pi^2/6$。}{For the square wave $f=-1$ on $(-\pi,0)$, $f=1$ on $(0,\pi)$, oddness gives $a_n=0$ and $b_n=4/(n\pi)$ for odd $n$, so $f\sim\frac4\pi\sum_{k\ge0}\frac{\sin((2k+1)x)}{2k+1}$; at $x=\pi/2$ this is Leibniz's formula. Parseval gives $\sum_{k\ge0}1/(2k+1)^2=\pi^2/8$, hence $\sum1/n^2=\pi^2/6$.}
\BOTH{\textbf{Gibbs 现象.} 部分和在跳跃点附近超调约 $8.95\%$（最大 $1.17898$），超调不随 $N$ 消失、只向跳跃点靠近，故部分和逐点收敛但在跳跃点附近不一致收敛。}{\textbf{Gibbs phenomenon.} Near a jump the partial sums overshoot by about $8.95\%$ of the jump, and the overshoot does not vanish as $N$ grows, only moving closer to the jump; the partial sums converge pointwise but not uniformly near the jump.}
\REMARK{应用：Fourier 系数给频谱、滤波即按频率修改系数；热传导 $u_t=a^2u_{xx}$ 分离变量得 $u=\sum b_n\sin(n\pi x/L)e^{-a^2(n\pi/L)^2t}$，高频指数衰减解释热平滑；FFT 在 $O(N\log N)$ 内近似计算系数。}{Applications: Fourier coefficients give the spectrum and filtering modifies them frequency by frequency; for heat conduction $u_t=a^2u_{xx}$ separation gives $u=\sum b_n\sin(n\pi x/L)e^{-a^2(n\pi/L)^2t}$, whose exponential decay explains smoothing; the FFT approximates coefficients in $O(N\log N)$.}
\ASK{Taylor 近似与 Fourier 近似有何不同？}{How do Taylor and Fourier approximations differ?}{Taylor 用一点的高阶导数描述局部行为，误差由高阶导数控制；Fourier 用整个周期上的正交积分系数描述全局周期结构，误差由系数衰减控制。前者局部，后者全局。}{Taylor uses higher derivatives at one point to describe local behaviour, with error controlled by higher derivatives; Fourier uses orthogonal integral coefficients over a whole period to describe global periodic structure, with error controlled by coefficient decay. The former is local, the latter global.}

\section*{中英术语表 | Glossary}
\begingroup\small
\setlength{\tabcolsep}{4pt}
\renewcommand{\arraystretch}{1.12}
\begin{tabular}{p{2.9cm}p{4.4cm}p{2.9cm}p{4.4cm}}
\hline
\textbf{中文} & \textbf{English} & \textbf{中文} & \textbf{English}\\
\hline
上确界 & supremum & 完备性 & completeness\\
单调有界定理 & monotone convergence & 闭区间套 & nested intervals\\
Cauchy 准则 & Cauchy criterion & 一致连续 & uniform continuity\\
连续 & continuity & 可导 & differentiable\\
微分 & differential & 中值定理 & mean-value theorem\\
Taylor 公式 & Taylor's formula & 原函数 & antiderivative\\
可积 & integrable & 微积分基本定理 & fundamental theorem of calculus\\
反常积分 & improper integral & 数项级数 & numerical series\\
比较判别法 & comparison test & 比值判别法 & ratio test\\
交错级数 & alternating series & 绝对收敛 & absolute convergence\\
条件收敛 & conditional convergence & 偏导数 & partial derivative\\
全微分 & total differential & 梯度 & gradient\\
链式法则 & chain rule & Jacobian 行列式 & Jacobian determinant\\
隐函数定理 & implicit function theorem & Lagrange 乘子 & Lagrange multiplier\\
重积分 & multiple integral & Fubini 定理 & Fubini's theorem\\
曲线积分 & line integral & 曲面积分 & surface integral\\
Green 公式 & Green's theorem & Stokes 公式 & Stokes' theorem\\
散度 & divergence & 旋度 & curl\\
保守场 & conservative field & Fourier 系数 & Fourier coefficients\\
Parseval 等式 & Parseval's identity & Gibbs 现象 & Gibbs phenomenon\\
含参积分 & integral with a parameter & 一致收敛 & uniform convergence\\
\hline
\end{tabular}
\endgroup

\section*{参考文献 | References}
\begin{enumerate}[label={[\arabic*]},leftmargin=2.4em]
\item 李逸. 《基本分析讲义》（东南大学理班讲义）. 第二章“极限理论 I: 数列极限”、第三章“极限理论 II: 函数极限”、第四章“导数理论”、第五章“积分理论”、第六章“级数理论”、第十一章“多变量极限理论”、第十二章“多变量导数理论”、第十三章“多变量积分理论”、第十六章“Fourier 级数”. 本讲义的定义、定理条件与例子主要依据这些章节核对. [Li Yi, \emph{Basic Analysis Lecture Notes} (Southeast University science-class notes), Chapters 2, 3, 4, 5, 6, 11, 12, 13 and 16.]
\item 华东师范大学数学系. 《数学分析》（第四版）. 北京: 高等教育出版社. 上册第七、十二、十三、十四、十五、十六章（实数完备性、极限、连续、导数、积分、级数、反常积分、Fourier 级数）与下册第十六、十七、十八、二十、二十一、二十二章（多元极限与连续、多元微分、隐函数、曲线积分、重积分、曲面积分）. [Department of Mathematics, East China Normal University, \emph{Mathematical Analysis}, 4th ed., Higher Education Press; chapters on completeness, limits, continuity, differentiation, integration, series and Fourier series (vol. 1), and on multivariable limits, differentiation, implicit functions, line and surface integrals, multiple integrals (vol. 2).]
\item W. Rudin. \emph{Principles of Mathematical Analysis}, 3rd ed. New York: McGraw--Hill, 1976. 实数构造、极限、连续性、微分、Riemann--Stieltjes 积分、函数序列与级数、多变量微分与积分的标准英文参考. [Standard English reference for the construction of the reals, limits, continuity, differentiation, integration, sequences and series of functions, and multivariable calculus.]
\end{enumerate}

\section*{范围与复习说明 | Scope note}
\BOTH{依据李逸《基本分析讲义》第二章至第六章、第十一章至第十三章与第十六章的指定内容；带星号的小节（微分形式、Fourier 级数一致收敛的详细理论、Lebesgue 可积判别等）不纳入；含参积分一节按第三册考纲补充。复习优先：每个定理的条件、结论、证明思路与一个反例；每个公式的适用条件；中英术语的准确说法。三册重复的数值例题已合并删除，只保留每章一两个短例或反例。}{Based on the designated content of Chapters 2--6, 11--13 and 16 of Li Yi's notes; starred subsections (differential forms, the detailed theory of uniform convergence of Fourier series, the Lebesgue criterion) are excluded. The section on integrals depending on a parameter supplements the third-semester syllabus. Review priorities: the hypotheses, conclusion, proof idea and one counterexample of each theorem; the conditions of each formula; and accurate bilingual terminology. Duplicated numerical examples across the three volumes have been removed, keeping only one or two short examples or counterexamples per chapter.}
\end{document}
